\documentclass[11pt,twoside]{article}

\usepackage[margin=1in]{geometry}
\usepackage{amsmath,amssymb,amsthm,mathtools}
\usepackage[inline,shortlabels]{enumitem}
\usepackage{booktabs,array}
\usepackage{xcolor}
\usepackage{microtype}
\usepackage{needspace}

% Revision date: maintained by scripts/build_lecture_notes.py. Keep this command.
\newcommand{\lecturerevised}{2026-10-25 21:25 UTC-06:00}
\usepackage{eso-pic}
\AddToShipoutPictureBG{%
  \AtPageLowerLeft{%
    \put(\LenToUnit{0.5\paperwidth},\LenToUnit{0.5in}){%
      \makebox(0,0){\normalfont\footnotesize\color{black!60}Last revised: \lecturerevised}%
    }%
  }%
}

\usepackage[square,authoryear]{natbib}
\usepackage[colorlinks=true,linkcolor=blue!60!black,
  urlcolor=blue!60!black,citecolor=blue!60!black]{hyperref}
\usepackage[nameinlink,capitalize]{cleveref}
\crefname{exercise}{Exercise}{Exercises}
\Crefname{exercise}{Exercise}{Exercises}
\hypersetup{pdftitle={Lecture 12: Accelerated and adaptive first-order methods},
  pdfauthor={Csaba Szepesvari}}

\usepackage{../exercises/course-exercises}

\usepackage{enotez}
\DeclareInstance{enotez-list}{lecture}{list}{format=\normalsize,number={#1.}}
\setenotez{list-name=Endnotes,list-style=lecture,backref=true}
\crefname{endnote}{Endnote}{Endnotes}
\Crefname{endnote}{Endnote}{Endnotes}

\definecolor{courseblue}{HTML}{17365D}
\newcommand{\E}{\mathbb E}
\newcommand{\I}{\mathbb I}
\renewcommand{\P}{\mathbb P}
\newcommand{\R}{\mathbb R}
\newcommand{\argmin}{\operatorname*{arg\,min}}
\newcommand{\diag}{\operatorname{diag}}
\newcommand{\suppmark}{\texorpdfstring{\textsuperscript{*}}{*}}
\makeatletter
\@for\@letter:=A,B,C,D,E,F,G,H,I,J,K,L,M,N,O,P,Q,R,S,T,U,V,W,X,Y,Z\do{%
  \expandafter\edef\csname c\@letter\endcsname{\noexpand\mathcal{\@letter}}}
\makeatother

\newtheorem{theorem}{Theorem}[section]
\newtheorem{proposition}[theorem]{Proposition}
\newtheorem{corollary}[theorem]{Corollary}
\newtheorem{lemma}[theorem]{Lemma}
\theoremstyle{definition}
\newtheorem{definition}[theorem]{Definition}
\newtheorem{example}[theorem]{Example}
\theoremstyle{remark}
\newtheorem{remark}[theorem]{Remark}

\renewcommand{\thepage}{12--\arabic{page}}
\renewcommand{\thesection}{12.\arabic{section}}
\renewcommand{\theequation}{12.\arabic{equation}}
\pagestyle{myheadings}
\markboth{Lecture 12: Accelerated and adaptive first-order methods}
{Lecture 12: Accelerated and adaptive first-order methods}

\begin{document}

\begin{center}
\fcolorbox{courseblue}{white}{%
\begin{minipage}{0.92\textwidth}
\vspace{1mm}
\textbf{CMPUT 654: Mathematical Foundations of Modern AI Systems}
\hfill Fall 2026

\vspace{3mm}
\begin{center}
{\Large\bfseries Lecture 12: Accelerated and adaptive first-order methods}

\vspace{1mm}
October 13, 2026
\end{center}

\vspace{2mm}
\textit{Lecturer: Csaba Szepesv\'ari}
\hfill
\textit{Note prepared by: Csaba Szepesv\'ari}
\vspace{1mm}
\end{minipage}}
\end{center}

\noindent\textbf{Reading guide.}
Lecture~11 showed that a $\beta$-smooth objective decreases under gradient
descent when the stepsize satisfies $\eta\le1/\beta$. Under a $\mu$-PL
inequality, the same update reduces the objective gap geometrically, and the
number of iterations needed for a fixed relative reduction is proportional
to $\beta/\mu$. Two questions remain.
First, an implementation may not know $\beta$. Second, even the best safe
scalar step can be slow when different directions have very different
curvature.

The two questions concern different limitations of gradient descent. An
unknown $\beta$ prevents the algorithm from choosing a certified scalar
stepsize. A large ratio $\beta/\mu$ means that every safe scalar stepsize is
slow in some directions. Backtracking addresses the first limitation by
testing proposed finite steps. Acceleration and preconditioning address the
second limitation by changing how information from different iterations or
different parameter directions is used. Keeping stepsize adaptation separate
from conditioning improvement is important: a method can choose a safe
stepsize automatically and still
converge slowly on an ill-conditioned problem.

The first part of this lecture develops backtracking line search, which tests
the finite-step inequality needed by the proof and therefore avoids advance
knowledge of $\beta$. The second part explains the condition-number barrier
for gradient descent and proves the main acceleration result: for smooth,
strongly convex objectives, the number of gradient evaluations needed to
reduce the objective gap by a fixed relative amount improves from order
$\kappa$ to order $\sqrt\kappa$. The third part separates scalar stepsize
adaptation from adaptation of the update geometry. It first studies
preconditioning for a fixed smooth objective and then changes the information
available to the update: AdaGrad and Adam adapt to the subgradients or
stochastic gradients observed during the run. AdamW adds a separate choice
about weight decay.

The principal results for a first reading are the backtracking guarantee, the
linear acceleration theorem and the corresponding lower-bound comparison, the
preconditioning calculation, and the distinction between scalar online
adaptation, diagonal AdaGrad, and Adam. Material marked with $\suppmark$ may
be skipped on a first reading. The prerequisites are the smoothness,
convexity, strong-convexity, and PL results from Lecture~11, together with
eigenvalues of symmetric matrices.

\section{Review: Gradient descent with smooth objectives}
\label{sec:tuning-problem}

This section recalls why the usual gradient-descent stepsize depends on the
smoothness constant. The argument also separates the condition used by the
convergence proof---a one-step decrease inequality---from one sufficient
certificate for that condition---a known global smoothness constant.

Let $L:\R^d\to\R$ be differentiable and bounded below, and put
\[
 L^\star=\inf_{\theta\in\R^d}L(\theta).
\]
The function is $\beta$-smooth when
\begin{equation}
 \|\nabla L(u)-\nabla L(v)\|_2
 \le\beta\|u-v\|_2
 \qquad\text{for every }u,v\in\R^d.
\label{eq:smoothness}
\end{equation}
Smoothness converts first-order information at $u$ into the quadratic upper
bound
\begin{equation}
 L(v)\le L(u)
 +\langle\nabla L(u),v-u\rangle
 +\frac\beta2\|v-u\|_2^2.
\label{eq:quadratic-upper-bound}
\end{equation}
Applying \cref{eq:quadratic-upper-bound} with
$u=\theta$ and $v=\theta^+=\theta-\eta\nabla L(\theta)$ gives
\begin{equation}
 L(\theta^+)
 \le L(\theta)-\eta\left(1-\frac{\beta\eta}{2}\right)
                  \|\nabla L(\theta)\|_2^2.
\label{eq:smooth-gradient-decrease}
\end{equation}
The usual stepsize restriction now enters at one precise point. If
$\eta\le1/\beta$, then
\begin{equation}
 L(\theta^+)\le L(\theta)-\frac\eta2\|\nabla L(\theta)\|_2^2.
\label{eq:sufficient-decrease-eta}
\end{equation}

The role of $\beta$ is therefore quite specific. Knowledge of $\beta$
guarantees \cref{eq:sufficient-decrease-eta} at every parameter, for every
step with $\eta\le1/\beta$. The convergence proof, however, uses the decrease
inequality itself; the proof does not otherwise use the numerical value of
$\beta$.
Moreover, the argument gives only a sufficient restriction on the stepsize.
A larger step may still satisfy the same decrease inequality at the current
parameter. Section~\ref{sec:backtracking} uses the distinction between the
decrease inequality and the numerical value of $\beta$ to construct a
stepsize rule: propose a step, test the required decrease, and shorten the
step only when the test fails.

\section{Adapting to an unknown smoothness constant}
\label{sec:backtracking}

This section develops a stepsize rule for the common situation in which a
useful value of $\beta$ is not known. Sometimes the available analysis gives
only a conservative bound $B\ge\beta$. The stepsize $1/B$ is safe, but a
large ratio $B/\beta$ makes gradient descent take smaller steps, and possibly
more iterations, than the sharper bound would require. The first part of
\cref{ex:l12-backtracking-quadratic} asks you to quantify the additional
iteration cost on a quadratic objective. In other problems, no finite upper
bound on $\beta$ may be available to the practitioner.

The preceding decrease calculation suggests searching for a safe upper
estimate of $\beta$. At a parameter $\theta$, let $g=\nabla L(\theta)$ and
let $M>0$ be a tentative upper bound on $\beta$. The estimate $M$ defines the
trial point
\begin{equation}
 \theta^+(M)=\theta-\frac1M g.
\label{eq:trial-point}
\end{equation}
If $M$ is indeed an upper bound, so that $M\ge\beta$, then
\cref{eq:smooth-gradient-decrease} with $\eta=1/M$ gives
\[
 L(\theta^+(M))
 \le L(\theta)
 -\frac1M\left(1-\frac{\beta}{2M}\right)\|g\|_2^2
 \le L(\theta)-\frac1{2M}\|g\|_2^2.
\]
Thus $M\ge\beta$ is a universally safe range: every estimate in that range
guarantees a decrease of at least $\|g\|_2^2/(2M)$. The factor $1/2$ is the
largest constant certified uniformly over all $\beta\le M$, because the worst
case $\beta=M$ makes the preceding coefficient exactly $1/2$.

Although the value of $\beta$ is unknown, the promised decrease can be tested
using function values. For a tentative estimate $M$, accept the trial point
when
\begin{equation}
 L(\theta^+(M))
 \le L(\theta)-\frac1{2M}\|g\|_2^2.
\label{eq:backtracking-test}
\end{equation}
Every quantity in \cref{eq:backtracking-test} is available: the gradient $g$
determines $\theta^+(M)$, and one objective evaluation gives
$L(\theta^+(M))$. If the test fails, increase the estimate geometrically by
replacing $M$ with $\rho M$ for a fixed $\rho>1$, form the new trial point,
and test again. Repeating these operations is called \emph{backtracking line
search}. The search variable is the curvature estimate $M$; after $M$ is
accepted, the gradient step uses the corresponding stepsize $1/M$.

The search must terminate because every $M\ge\beta$ passes
\cref{eq:backtracking-test}. Starting from $M_{\rm start}$, after $k$ failures
the estimate is $\rho^kM_{\rm start}$, so a safe estimate is reached after
only logarithmically many trials. If $M_{\rm start}<\beta$, the first estimate
that reaches the safe range is smaller than $\rho\beta$. Geometric increases
therefore find an estimate within a factor $\rho$ of the threshold $\beta$;
an additive increase would depend on the numerical scale of $\beta$.

The name \emph{backtracking} describes the movement of the trial point. Each
increase of $M$ shortens the step and moves $\theta^+(M)$ back toward
$\theta$ along the negative-gradient direction.

The acceptance test is local and can succeed before $M$ reaches $\beta$. The
constant $\beta$ controls the quadratic upper bound for every pair of
parameters, whereas \cref{eq:backtracking-test} concerns only $\theta$ and the
proposed successor $\theta^+(M)$. If the test accepts a value $M<\beta$, the
algorithm takes a step larger than $1/\beta$ without requiring the same
curvature bound to hold elsewhere.

\paragraph{Backtracking gradient descent.}
The complete method uses an initial estimate $M_{\mathrm{init}}>0$ and a
multiplier $\rho>1$. At iteration $t$:
\begin{enumerate}[(1),leftmargin=*]
\item compute $g_t=\nabla L(\theta_t)$;
\item set $M=M_{\mathrm{init}}$ when $t=0$, and set $M=M_{t-1}$ when
$t\ge1$;
\item form $\theta^+(M)=\theta_t-g_t/M$ and test whether
\[
 L(\theta^+(M))
 \le L(\theta_t)-\frac1{2M}\|g_t\|_2^2.
\]
If the test fails, replace $M$ by $\rho M$ and repeat this item;
\item set $M_t=M$ and $\theta_{t+1}=\theta^+(M_t)$.
\end{enumerate}
Rejected trials do not require additional gradient evaluations. The current
parameter has not changed, so the method reuses $g_t$ and pays only for the
objective values at the new trial points.

At every iteration $t\ge1$, the algorithm begins with $M_{t-1}$.
Consequently, $(M_t)$ is nondecreasing, and every rejected trial causes a
permanent geometric increase in the estimate. The monotone rule gives the
simple one-time bound on rejected trials proved below. The limitation is that
the method may miss a smaller locally safe estimate, and hence a larger safe
stepsize, at a later iteration.

\begin{example}[The safe range is exact on a scalar quadratic]
\label{ex:scalar-backtracking}
The smoothness argument proves that every $M\ge\beta$ passes the acceptance
test. The scalar quadratic shows that this sufficient condition is also
necessary in the worst case. Let $L(\theta)=\beta\theta^2/2$. Then
\[
 \theta^+(M)=\left(1-\frac\beta M\right)\theta.
\]
For $\theta\ne0$, substituting into \cref{eq:backtracking-test} and dividing
by $\beta\theta^2/2$ gives the exact equivalence
\begin{equation}
 \left(1-\frac\beta M\right)^2
 \le 1-\frac\beta M
 \quad\Longleftrightarrow\quad
 M\ge\beta.
 \label{eq:scalar-backtracking-threshold}
\end{equation}
Thus the threshold $\beta$ is
not an artifact of the proof: on this objective, every smaller estimate
fails. More general objectives can still accept $M<\beta$ when the curvature
encountered by the current step is smaller than the global worst case.
\Cref{ex:l12-backtracking-quadratic} asks you to analyze both the exact
quadratic threshold and the possible gain from local curvature.
\end{example}

\begin{proposition}[Termination and a bound on $M_t$]
\label{prop:backtracking-termination}
Suppose that $L$ is $\beta$-smooth. Run the backtracking gradient descent
defined above with $M_{\mathrm{init}}>0$ and $\rho>1$. Every search terminates
after finitely many trials, and, for every $t\ge0$,
\begin{equation}
 M_t\le\max\{M_{\mathrm{init}},\rho\beta\}.
 \label{eq:accepted-scale}
\end{equation}
\end{proposition}

\begin{proof}
The calculation preceding \cref{eq:backtracking-test} shows that every
$M\ge\beta$ passes the test, so each search terminates. We prove
\cref{eq:accepted-scale} by induction on $t$. At $t=0$, the algorithm first
tests $M_{\mathrm{init}}$. If the first trial passes, then
$M_0=M_{\mathrm{init}}$. Otherwise, the algorithm multiplies the trial value
by $\rho$ until $M_0$ passes. The immediately preceding value $M_0/\rho$ was
rejected and must therefore satisfy $M_0/\rho<\beta$. Hence either
$M_0=M_{\mathrm{init}}$ or $M_0<\rho\beta$, which proves
\cref{eq:accepted-scale} at $t=0$.

Now suppose the bound holds at $t-1$. If iteration $t$ accepts its first
trial, then $M_t=M_{t-1}$ and the induction hypothesis applies. If iteration
$t$ rejects at least one trial, then its last rejected value is
$M_t/\rho<\beta$, so $M_t<\rho\beta$. The bound follows in both cases.
\end{proof}

Proposition~\ref{prop:backtracking-termination} also bounds the extra
objective evaluations used by the search. Across all iterations, the number
of rejected trial points is
at most
\begin{equation}
 N_{\rm reject}
 \le
 \max\left\{0,
   \left\lceil\log_\rho\frac{\beta}{M_{\mathrm{init}}}\right\rceil
 \right\}.
 \label{eq:backtracking-rejections}
\end{equation}
Every rejection multiplies the current estimate by $\rho$, the sequence
$(M_t)$ is nondecreasing, and every estimate at least $\beta$ passes the test.
Hence $T$
iterations require at most $T+N_{\rm reject}$ objective evaluations at trial
points, in addition to the $T$ gradient evaluations.

\subsection{The deterministic price is a constant}

Proposition~\ref{prop:backtracking-termination} and
\cref{eq:backtracking-rejections} bound the extra objective evaluations used
to choose the steps. The next theorem bounds the number of gradient
iterations: under the PL inequality, the gradient-iteration bound needed to
reach a prescribed objective gap is at most a factor $\rho$ larger than the
corresponding known-$\beta$ bound.

\begin{theorem}[Backtracking under PL]
\label{thm:backtracking-pl}
Suppose that $L$ is $\beta$-smooth and satisfies
\begin{equation}
 \frac12\|\nabla L(\theta)\|_2^2
 \ge\mu(L(\theta)-L^\star)
 \qquad\text{for every }\theta\in\R^d.
 \label{eq:pl-condition}
\end{equation}
Run the backtracking gradient descent defined above with multiplier $\rho>1$
and initial estimate $M_{\mathrm{init}}>0$. Put
$\overline M=\max\{M_{\mathrm{init}},\rho\beta\}$. Then
\begin{equation}
 L(\theta_t)-L^\star
 \le\left(1-\frac\mu{\overline M}\right)^t
      \bigl(L(\theta_0)-L^\star\bigr).
 \label{eq:backtracking-pl-rate}
\end{equation}
\end{theorem}

\begin{proof}
Acceptance and \cref{eq:pl-condition} give
\begin{align*}
 L(\theta_{t+1})-L^\star
 &\le L(\theta_t)-L^\star
      -\frac1{2M_t}\|\nabla L(\theta_t)\|_2^2\\
 &\le\left(1-\frac\mu{M_t}\right)
      \bigl(L(\theta_t)-L^\star\bigr).
\end{align*}
By \cref{prop:backtracking-termination}, $M_t\le\overline M$. Replacing
$\mu/M_t$ by the smaller number $\mu/\overline M$ and iterating proves the
claim.
\end{proof}

With the fixed stepsize $1/\beta$, the objective gap satisfies
\[
 L(\theta_{t+1})-L^\star
 \le\left(1-\frac\mu\beta\right)
       \bigl(L(\theta_t)-L^\star\bigr).
\]
If $M_{\mathrm{init}}\le\beta$, backtracking instead guarantees the one-step
objective-gap multiplier $1-\mu/(\rho\beta)$ and adds the trial objective
evaluations counted above.
The number of gradient iterations remains
\[
 O\left(\frac\beta\mu
 \log\frac{L(\theta_0)-L^\star}{\epsilon}\right).
\]
Thus adaptation to an unknown deterministic smoothness constant changes the
iteration bound only by the constant factor $\rho$. Backtracking may do
better than the worst-case factor-$\rho$ comparison when the values $M_t$
that pass the local test are smaller than the global constant $\beta$.

The conclusion relies on access to exact objective values at the trial
points. A noisy minibatch loss does not generally certify the same inequality
for the full empirical objective or the population objective. Adaptation in
the stochastic setting therefore requires a different argument.

\section{The condition number and gradient descent}
\label{sec:conditioning}

This section identifies the limitation that remains after backtracking has
found a safe scalar step. One scalar must accommodate every parameter
direction, so a step that is stable in a high-curvature direction can make
little progress in a low-curvature direction. A strongly convex quadratic
makes this condition-number obstruction exact and prepares the acceleration
analysis in the next section.

Consider
\begin{equation}
 L(\theta)=\frac12\theta^\top H\theta-b^\top\theta,
 \qquad
 0<\mu I\preceq H\preceq\beta I.
 \label{eq:quadratic-objective}
\end{equation}
The quadratic objective has the unique minimizer
$\theta^\star=H^{-1}b$. If
$e_t=\theta_t-\theta^\star$, gradient descent gives
\begin{equation}
 e_{t+1}=(I-\eta H)e_t.
 \label{eq:quadratic-error-recursion}
\end{equation}
After diagonalizing $H=Q\diag(\lambda_1,\ldots,\lambda_d)Q^\top$, every
eigencomponent evolves as
\begin{equation}
 \langle q_j,e_t\rangle
 =(1-\eta\lambda_j)^t\langle q_j,e_0\rangle.
 \label{eq:eigencomponent-gd}
\end{equation}

For this quadratic, the sharp choices of the smoothness constant and the
strong-convexity or PL constant are $\beta=\lambda_{\max}(H)$ and
$\mu=\lambda_{\min}(H)$. With these choices, the condition number defined in
Lecture~11 equals the spectral ratio
\begin{equation}
 \kappa=\frac\beta\mu
 =\frac{\lambda_{\max}(H)}{\lambda_{\min}(H)}.
 \label{eq:condition-number}
\end{equation}
If $\beta$ and $\mu$ are instead loose bounds in
$\mu I\preceq H\preceq\beta I$, then $\beta/\mu$ is an upper bound on the
spectral ratio.
A step small enough to be stable in the largest-curvature direction changes
a smallest-curvature component by only a fraction of order $1/\kappa$ per
iteration.

The two-dimensional diagonal case makes the obstruction explicit. For
\[
 L(\theta)=\frac12\bigl(\mu\theta_1^2+\beta\theta_2^2\bigr),
\]
the step $\eta=1/\beta$ eliminates the second coordinate in one iteration,
whereas the first coordinate obeys
\[
 \theta_{t+1,1}=\left(1-\frac1\kappa\right)\theta_{t,1}.
\]
The same scalar stepsize must accommodate both directions. Backtracking can
discover that $1/\beta$ is safe, but discovering the safe value does not
change the recursion
$\theta_{t+1,1}=(1-1/\kappa)\theta_{t,1}$ in the flatter direction.

\begin{proposition}[Best constant step on a quadratic]
\label{prop:best-constant-step}
For \cref{eq:quadratic-objective}, define
\begin{equation}
 q(\eta)=\max_{\lambda\in[\mu,\beta]}|1-\eta\lambda|.
 \label{eq:quadratic-uniform-contraction}
\end{equation}
For every $H$ satisfying $\mu I\preceq H\preceq\beta I$, gradient descent
satisfies
\begin{equation}
 \|e_{t+1}\|_2\le q(\eta)\|e_t\|_2.
 \label{eq:quadratic-error-one-step}
\end{equation}
The constant step
\begin{equation}
 \eta_{\rm opt}=\frac{2}{\beta+\mu}
 \label{eq:optimal-quadratic-step}
\end{equation}
minimizes $q(\eta)$ over $\eta>0$:
\begin{equation}
 \min_{\eta>0}q(\eta)
 =q(\eta_{\rm opt})
 =\frac{\kappa-1}{\kappa+1}.
 \label{eq:gd-spectral-factor}
\end{equation}
Consequently,
\begin{equation}
 \|e_t\|_2
 \le\left(\frac{\kappa-1}{\kappa+1}\right)^t\|e_0\|_2.
 \label{eq:quadratic-error-t-step}
\end{equation}
If $e_0\ne0$ and $He_0=\lambda e_0$ for some
$\lambda\in\{\mu,\beta\}$, then equality holds in
\cref{eq:quadratic-error-t-step}. Hence, when $\kappa>1$, for every fixed
$c\in(0,1)$ and every such pair $(H,e_0)$,
\begin{equation}
 \min\{t\ge0:\|e_t\|_2\le c\|e_0\|_2\}
 =\left\lceil
   \frac{\log(1/c)}{\log((\kappa+1)/(\kappa-1))}
  \right\rceil
 =\Theta(\kappa)
 \qquad(\kappa\to\infty).
 \label{eq:quadratic-fixed-reduction}
\end{equation}
\end{proposition}

\begin{proof}
Diagonalizing $H$ in \cref{eq:quadratic-error-recursion} gives
\[
 \|e_{t+1}\|_2
 \le \max_j|1-\eta\lambda_j|\,\|e_t\|_2
 \le q(\eta)\|e_t\|_2,
\]
which proves \cref{eq:quadratic-error-one-step}. For fixed $\eta>0$, the
function $\lambda\mapsto |1-\eta\lambda|$ is convex, so its maximum on
$[\mu,\beta]$ occurs at $\lambda=\mu$ or $\lambda=\beta$. Expanding the two
absolute values gives four affine functions. The inequalities
$1-\eta\beta\le 1-\eta\mu$ and
$\eta\mu-1\le\eta\beta-1$ remove the two dominated functions, so
$q(\eta)=\max\{1-\eta\mu,\eta\beta-1\}$. The first function decreases with
$\eta$, the second increases with $\eta$, and they are equal exactly at
$\eta=\eta_{\rm opt}=2/(\beta+\mu)$. Consequently,
\[
 q(\eta)=
 \begin{cases}
  1-\eta\mu,&0<\eta\le\eta_{\rm opt},\\
  \eta\beta-1,&\eta\ge\eta_{\rm opt}.
 \end{cases}
\]
The maximum therefore decreases up to $\eta_{\rm opt}$ and increases after
$\eta_{\rm opt}$, so its minimum occurs at $\eta_{\rm opt}$. Substitution
gives \cref{eq:gd-spectral-factor}. Iterating
\cref{eq:quadratic-error-one-step} proves
\cref{eq:quadratic-error-t-step}. If $He_0=\lambda e_0$ for
$\lambda\in\{\mu,\beta\}$, then $e_t=(1-\eta_{\rm opt}\lambda)^t e_0$.
For $\lambda=\mu$, the multiplier is
$(\kappa-1)/(\kappa+1)$; for $\lambda=\beta$, the multiplier is
$-(\kappa-1)/(\kappa+1)$. Hence
$\|e_t\|_2=((\kappa-1)/(\kappa+1))^t\|e_0\|_2$ in both cases. The sign
alternates when $\lambda=\beta$, but the norm is unchanged. Thus equality
holds in
\cref{eq:quadratic-error-t-step}. Solving
$q(\eta_{\rm opt})^t\le c$ gives
the exact iteration count in \cref{eq:quadratic-fixed-reduction}. The
$\Theta(\kappa)$ relation follows from
$\log((\kappa+1)/(\kappa-1))=2/\kappa+O(1/\kappa^2)$.
\end{proof}

The error after $t$ steps has the form
\[
 e_t=p_t(H)e_0,
 \qquad p_t(\lambda)=(1-\eta\lambda)^t.
\]
A method that combines past iterates can generate a degree-$t$ polynomial
with a smaller maximum magnitude on $[\mu,\beta]$. Acceleration chooses such
a polynomial without computing an eigendecomposition.

The polynomial formulation also explains why acceleration concerns more than
a larger stepsize. Repeated gradient descent applies the same linear factor
to every eigenvalue at every iteration. A method with memory can choose the
coefficients of a degree-$t$ polynomial to reduce
$\max_{\lambda\in[\mu,\beta]}|p_t(\lambda)|$. The iterates implement the
recurrence for the chosen polynomial using gradients and a small amount of
stored state.

\section{Reducing the condition-number dependence with acceleration}
\label{sec:acceleration}

This section asks whether a first-order method can make faster progress in
low-curvature directions without taking an unstable gradient step in
high-curvature directions. The backtracking procedure used smoothness, and its
linear-rate theorem added the PL inequality without convexity. The direct
acceleration argument below narrows the setting.
Throughout the derivation, $L$ is $\beta$-smooth and $\mu$-strongly convex,
and $\theta^\star$ is its unique minimizer. Acceleration retains the preceding
displacement and reduces the iteration-count dependence from $\kappa$ to
$\sqrt\kappa$.

The design begins by deciding what the method may store and what quantity the
analysis should decrease. Gradient descent stores one parameter vector. We
permit one additional parameter-sized vector $v_t\in\R^d$ and treat
$(\theta_t,v_t)$ as the state. Thus the method uses twice the parameter
memory of gradient descent, apart from the gradient and temporary vectors.
For an undetermined coefficient $\nu>0$, consider the Lyapunov function
\begin{equation}
 V(\theta,v)
 =L(\theta)-L^\star+\frac\nu2\|v-\theta^\star\|_2^2.
 \label{eq:accelerated-design-potential}
\end{equation}
The objective gap is the quantity that we want to control. The second term
allows the analysis to record progress that is stored in the auxiliary
state. Since
\begin{equation}
 L(\theta)-L^\star\le V(\theta,v)
 \qquad\text{for every }\theta,v\in\R^d,
 \label{eq:accelerated-potential-dominates-gap}
\end{equation}
a geometric bound on $V$ immediately gives the same bound on the
objective gap. We therefore seek an update and a number $\tau\in(0,1]$ such
that
\begin{equation}
 V(\theta_{t+1},v_{t+1})
 \le(1-\tau)V(\theta_t,v_t).
 \label{eq:accelerated-design-target}
\end{equation}

We search within a four-parameter update family. Requiring the intermediate
point to equal the common stored vector when $\theta_t=v_t$ suggests an affine
combination of $\theta_t$ and $v_t$. The parameter update takes a gradient step
from that point, and the auxiliary-state update combines the two state vectors
affinely and adds a gradient correction. A general linear two-state update
would introduce more coefficients. We use the restricted family below with
$b\in(0,1]$, so the state part of the $v_t$ update is a convex combination.
The verification will determine $a,b,c,\eta$, and the potential coefficient
$\nu$:
\begin{align}
 y_t&=(1-a)\theta_t+a v_t,\notag\\
 \theta_{t+1}&=y_t-\eta g_t,\qquad g_t=\nabla L(y_t),\notag\\
 v_{t+1}&=(1-b)v_t+b y_t+c g_t.
 \label{eq:accelerated-affine-search}
\end{align}

We now substitute the update directly into the next potential. Smoothness gives
\begin{align}
 V(\theta_{t+1},v_{t+1})
 &=L(\theta_{t+1})-L^\star
   +\frac\nu2\|v_{t+1}-\theta^\star\|_2^2\notag\\
 &\le L(y_t)-L^\star
 -\left(\eta-\frac{\beta\eta^2}{2}\right)\|g_t\|_2^2\notag\\
 &\quad+\frac\nu2\left\|(1-b)(v_t-\theta^\star)
       +b(y_t-\theta^\star)+c g_t\right\|_2^2.
 \label{eq:accelerated-design-initial-expansion}
\end{align}
The coefficient subtracted from $\|g_t\|_2^2$ is largest at $\eta=1/\beta$.
With that choice, expanding the squared norm gives
\begin{align}
 V(\theta_{t+1},v_{t+1})
 \le{}&L(y_t)-L^\star
 +\frac\nu2\left\|(1-b)(v_t-\theta^\star)
                    +b(y_t-\theta^\star)\right\|_2^2\notag\\
 &+\nu c\left\langle g_t,
          (1-b)(v_t-\theta^\star)+b(y_t-\theta^\star)\right\rangle
 +\left(\frac{\nu c^2}{2}-\frac1{2\beta}\right)\|g_t\|_2^2.
 \label{eq:accelerated-design-expanded}
\end{align}
The coefficient $b$ in the auxiliary-state update suggests trying
$\tau=b$ in the contraction target \cref{eq:accelerated-design-target}. The
state part of the update satisfies
\begin{equation}
 \left\|(1-b)(v_t-\theta^\star)+b(y_t-\theta^\star)\right\|_2^2
 =(1-b)\|v_t-\theta^\star\|_2^2+b\|y_t-\theta^\star\|_2^2
 -b(1-b)\|v_t-y_t\|_2^2.
 \label{eq:accelerated-convex-combination-identity}
\end{equation}
The identity assigns weight $1-b$ to the squared distance involving $v_t$ and
weight $b$ to the squared distance involving $y_t$. To use the same split for
the objective gap, fix $0<\nu\le\mu$ and apply strong convexity to its
$b$-weighted part:
\begin{equation}
 b\bigl(L(y_t)-L^\star\bigr)
 \le b\langle g_t,y_t-\theta^\star\rangle
 -\frac{b\nu}{2}\|y_t-\theta^\star\|_2^2.
 \label{eq:accelerated-design-strong-convexity}
\end{equation}
Combining
\cref{eq:accelerated-design-expanded,eq:accelerated-convex-combination-identity,eq:accelerated-design-strong-convexity}
then gives
\begin{align}
 V(\theta_{t+1},v_{t+1})
 \le{}&(1-b)\left(
 L(y_t)-L^\star+\frac\nu2\|v_t-\theta^\star\|_2^2
 -\frac{b\nu}{2}\|v_t-y_t\|_2^2\right)\notag\\
 &+b\langle g_t,y_t-\theta^\star\rangle
 +\nu c\left\langle g_t,
 (1-b)(v_t-\theta^\star)+b(y_t-\theta^\star)\right\rangle\notag\\
 &+\left(\frac{\nu c^2}{2}-\frac1{2\beta}\right)\|g_t\|_2^2.
 \label{eq:accelerated-design-after-strong-convexity}
\end{align}
The first line has the proposed factor $1-b$. The next two lines determine
the remaining coefficient choices.

At this point, the design has imposed
\begin{equation*}
 \eta=\frac1\beta,
 \qquad
 \tau=b,
 \qquad
 0<\nu\le\mu.
\end{equation*}
The values of $a,b,c$, and $\nu$ remain to be selected within these
restrictions. To eliminate $\theta^\star$ from the two inner products in
\cref{eq:accelerated-design-after-strong-convexity} and produce another factor
$1-b$, set
\begin{equation}
 \nu c=-b
 \qquad\Longleftrightarrow\qquad
 c=-\frac b\nu.
 \label{eq:accelerated-design-c-choice}
\end{equation}
Indeed, the two inner products then satisfy
\begin{align}
 &b\langle g_t,y_t-\theta^\star\rangle
 -b\left\langle g_t,
 (1-b)(v_t-\theta^\star)+b(y_t-\theta^\star)\right\rangle\notag\\
 &\qquad=-b(1-b)\langle g_t,v_t-y_t\rangle.
 \label{eq:accelerated-design-gradient-cancellation}
\end{align}
Substituting \cref{eq:accelerated-design-c-choice,eq:accelerated-design-gradient-cancellation}
into \cref{eq:accelerated-design-after-strong-convexity} gives
\begin{align}
 V(\theta_{t+1},v_{t+1})
 \le{}&(1-b)\left(
 L(y_t)-L^\star+\frac\nu2\|v_t-\theta^\star\|_2^2
 -\frac{b\nu}{2}\|v_t-y_t\|_2^2
 -b\langle g_t,v_t-y_t\rangle\right)\notag\\
 &+\left(\frac{b^2}{2\nu}-\frac1{2\beta}\right)\|g_t\|_2^2.
 \label{eq:accelerated-design-general-bound}
\end{align}

The squared-gradient term in
\cref{eq:accelerated-design-general-bound} is nonpositive when
$b^2\le\nu/\beta$. Under this condition,
\begin{equation}
 V(\theta_{t+1},v_{t+1})
 \le(1-b)\left(
 L(y_t)-L^\star+\frac\nu2\|v_t-\theta^\star\|_2^2
 -\frac{b\nu}{2}\|v_t-y_t\|_2^2
 -b\langle g_t,v_t-y_t\rangle
 \right).
 \label{eq:accelerated-design-before-convexity}
\end{equation}
Add and subtract $\langle g_t,\theta_t-y_t\rangle$ in
\cref{eq:accelerated-design-before-convexity}. Convexity gives
$L(y_t)+\langle g_t,\theta_t-y_t\rangle\le L(\theta_t)$, and hence
\begin{equation}
 V(\theta_{t+1},v_{t+1})
 \le(1-b)\left(V(\theta_t,v_t)-\frac{b\nu}{2}\|v_t-y_t\|_2^2
 +\left\langle g_t,b(y_t-v_t)-(\theta_t-y_t)\right\rangle\right).
 \label{eq:accelerated-design-after-convexity}
\end{equation}
Choose the intermediate point so that the last inner product vanishes:
\begin{equation}
 \theta_t-y_t=b(y_t-v_t).
 \label{eq:accelerated-design-displacement}
\end{equation}
For $y_t=(1-a)\theta_t+a v_t$,
\begin{align}
 \theta_t-y_t&=a(\theta_t-v_t),\notag\\
 b(y_t-v_t)&=b(1-a)(\theta_t-v_t).
\end{align}
Thus \cref{eq:accelerated-design-displacement} determines
\begin{equation}
 a=b(1-a)
 \qquad\Longleftrightarrow\qquad
 a=\frac{b}{1+b}.
 \label{eq:accelerated-coupling-coefficient}
\end{equation}
Using \cref{eq:accelerated-design-displacement} in
\cref{eq:accelerated-design-after-convexity} gives
\begin{equation}
 V(\theta_{t+1},v_{t+1})
 \le(1-b)\left(
 V(\theta_t,v_t)-\frac\nu{2b}\|\theta_t-y_t\|_2^2
 \right)
 \le(1-b)V(\theta_t,v_t).
 \label{eq:accelerated-design-contraction}
\end{equation}
The contraction factor is $1-b$, subject to $b^2\le\nu/\beta$ and
$0<\nu\le\mu$. The largest admissible $b$ is obtained from
\begin{equation}
 \nu=\mu,
 \qquad
 b=\tau=\sqrt{\frac\mu\beta}=\frac1{\sqrt\kappa}.
 \label{eq:accelerated-design-balance}
\end{equation}
This choice proves the target \cref{eq:accelerated-design-target}. The
completed coefficient choices give the coupled updates
\begin{align}
 y_t&=\frac{\theta_t+\tau v_t}{1+\tau},
 \label{eq:accelerated-design-coupling}\\
 \theta_{t+1}&=y_t-\frac1\beta\nabla L(y_t),
 \label{eq:accelerated-design-gradient-step}\\
 v_{t+1}&=(1-\tau)v_t+\tau\left(y_t-\frac1\mu\nabla L(y_t)\right).
 \label{eq:accelerated-design-center-step}
\end{align}

The auxiliary state $v_t$ was useful for choosing the coefficients and proving
the contraction. Eliminating $v_t$ gives an equivalent implementation that
stores the two most recent parameter iterates. Set $\theta_{-1}=\theta_0$ and
\begin{equation}
 \gamma=\frac{1-\tau}{1+\tau}
 =\frac{\sqrt\kappa-1}{\sqrt\kappa+1}.
 \label{eq:accelerated-momentum}
\end{equation}
The coupled update is equivalent to the two-step recursion
\begin{align}
 y_t&=\theta_t+\gamma(\theta_t-\theta_{t-1}),
 \label{eq:accelerated-extrapolation}\\
 \theta_{t+1}&=y_t-\frac1\beta\nabla L(y_t).
 \label{eq:accelerated-gradient-step}
\end{align}
This two-step recursion is the usual form of Nesterov's accelerated gradient
method. Exercise~\ref{ex:l12-nesterov-equivalence} asks you to verify the
equivalence and complete the proof of the accelerated rate.

The point $y_t$ extrapolates the displacement
$\theta_t-\theta_{t-1}$, and the second line takes an ordinary gradient step
from the extrapolated point. The update evaluates one gradient and stores two
parameter vectors per iteration.
The derivation illustrates the general strategy: specify the desired progress
measure, introduce state that can retain the missing information, and choose
the coefficients so that the available inequalities prove the desired
one-step change. Eliminating the auxiliary state then gives a simpler
implementation.

\begin{theorem}[Accelerated linear convergence]
\label{thm:accelerated-strongly-convex}
Suppose that $L$ is $\beta$-smooth and $\mu$-strongly convex, let
$\kappa=\beta/\mu$, and let $\theta^\star$ be its minimizer. Set $\gamma$ as
in \cref{eq:accelerated-momentum} and initialize $\theta_{-1}=\theta_0$. The
iterates defined by
\cref{eq:accelerated-extrapolation,eq:accelerated-gradient-step} satisfy
\begin{equation}
 L(\theta_t)-L^\star
 \le 2\left(1-\frac1{\sqrt\kappa}\right)^t
 \bigl(L(\theta_0)-L^\star\bigr).
 \label{eq:accelerated-linear-rate}
\end{equation}
Consequently, for
$0<\epsilon<2(L(\theta_0)-L^\star)$,
$L(\theta_t)-L^\star\le\epsilon$ whenever
\begin{equation}
 t\ge \sqrt\kappa
 \log\frac{2(L(\theta_0)-L^\star)}{\epsilon}.
 \label{eq:accelerated-strong-rate}
\end{equation}
\end{theorem}

The proof is completed in
\cref{ex:l12-nesterov-equivalence}. The preceding calculation supplies the
potential contraction \cref{eq:accelerated-design-contraction}; the exercise
eliminates $v_t$, bounds the initial potential using strong convexity, and
derives \cref{eq:accelerated-linear-rate,eq:accelerated-strong-rate}.

\paragraph{Quadratic interpretation.}
The quadratic calculation explains how the resulting method acts in different
curvature directions. Let $q$ be an eigenvector of $H$ with eigenvalue
$\lambda$, and write
\[
 e_t(\lambda)=\langle q,\theta_t-\theta^\star\rangle.
\]
For the quadratic objective in \cref{eq:quadratic-objective},
\cref{eq:accelerated-extrapolation,eq:accelerated-gradient-step} give
\begin{equation}
 e_{t+1}(\lambda)
 =\left(1-\frac{\lambda}{\beta}\right)
  \bigl((1+\gamma)e_t(\lambda)-\gamma e_{t-1}(\lambda)\bigr).
 \label{eq:accelerated-eigendirection}
\end{equation}
Gradient descent uses a first-order recurrence in each eigendirection;
accelerated gradient uses the second-order recurrence
\cref{eq:accelerated-eigendirection}. The calculation explains the role of
memory after the method has been derived. The related Chebyshev construction
in \cref{sec:chebyshev} solves the exact finite-degree minimax polynomial
problem for quadratics.

\paragraph{Nesterov and heavy-ball momentum.}
Nesterov's method evaluates the gradient at the extrapolated point $y_t$.
Polyak's heavy-ball method instead uses
\[
 \theta_{t+1}
 =\theta_t-\eta_{\rm HB}\nabla L(\theta_t)
  +\gamma_{\rm HB}(\theta_t-\theta_{t-1}).
\]
For a quadratic with spectrum in $[\mu,\beta]$, the quadratic-optimal fixed
parameters and the resulting worst-case asymptotic spectral factor are
\[
 \eta_{\rm HB}=\frac{4}{(\sqrt\beta+\sqrt\mu)^2},
 \qquad
 \gamma_{\rm HB}=r^2,
 \qquad
 r=\frac{\sqrt\kappa-1}{\sqrt\kappa+1}
   =1-\Theta\left(\frac1{\sqrt\kappa}\right).
\]
Thus heavy ball has accelerated $\sqrt\kappa$ dependence on strongly convex
quadratics. The quadratic guarantee is not uniform over all smooth, strongly
convex objectives: the same parameters can fail to converge on a
nonquadratic objective. By contrast,
\cref{thm:accelerated-strongly-convex} gives Nesterov's guarantee over the
entire smooth, strongly convex class. Section~\ref{sec:chebyshev} shows that
Chebyshev semi-iteration has the heavy-ball form with coefficients that vary
with $t$ and converge to $\eta_{\rm HB}$ and $\gamma_{\rm HB}$.

The gradient-evaluation comparison is
\begin{equation}
 \underbrace{O\left(\kappa
 \log\frac{L(\theta_0)-L^\star}{\epsilon}\right)}
 _{\text{gradient descent}}
 \qquad\text{and}\qquad
 \underbrace{O\left(\sqrt\kappa
 \log\frac{L(\theta_0)-L^\star}{\epsilon}\right)}
 _{\text{acceleration}}.
 \label{eq:gd-acceleration-comparison}
\end{equation}
Both methods require one gradient evaluation per iteration.
\Cref{ex:l12-accelerated-restart} asks you to derive the same square-root
dependence by restarting a smooth convex accelerated method.

The reduction from $\kappa$ to $\sqrt\kappa$ in the worst-case number of
first-order oracle calls is optimal in the standard first-order oracle model.
In this model, the method queries a parameter $\theta$ and receives
$L(\theta)$ and $\nabla L(\theta)$; each new query may depend on the preceding
queries and replies. The lower bounds also explain why the assumptions must be
stated carefully.

\begin{proposition}[Relevant first-order lower bounds]
\label{prop:acceleration-lower-bounds}
Fix $\kappa>1$ and consider dimensions large enough relative to the number of
oracle calls. Let $N_\epsilon$ denote the number of oracle calls required to
produce $\theta$ satisfying
\begin{equation}
 L(\theta)-L^\star
 \le\epsilon\bigl(L(\theta_0)-L^\star\bigr).
 \label{eq:relative-gap-target}
\end{equation}
\begin{enumerate}[(a),leftmargin=*]
\item For every deterministic first-order method, there is a smooth, strongly
convex quadratic instance on which
\begin{equation}
 N_\epsilon
 =\Omega\left(\sqrt\kappa\log\frac1\epsilon\right).
 \label{eq:strongly-convex-lower-bound}
\end{equation}
\item For the larger class of smooth objectives satisfying the PL inequality,
without a convexity assumption, every deterministic first-order method has an
instance on which
\begin{equation}
 N_\epsilon
 =\Omega\left(\kappa\log\frac1\epsilon\right).
 \label{eq:pl-lower-bound}
\end{equation}
\end{enumerate}
\end{proposition}

Part (a) shows that the square-root dependence on $\kappa$ is optimal over
the smooth, strongly convex class. Part (b) concerns the general smooth PL
class, which includes nonconvex objectives, and rules out a general
$O(\bigl(\sqrt\kappa\log(1/\epsilon)\bigr)$ guarantee over the full smooth PL
class. Smooth convex PL objectives form an intermediate class: the restart
construction below attains the square-root dependence for that class.

The smooth convex $O(1/t^2)$ result is a different acceleration theorem. The
main result here is the linear convergence bound
\cref{eq:accelerated-linear-rate}; the smooth convex result and a restart
argument are summarized in \cref{sec:restart-optional}.

\subsection{Accelerating when the condition number is unknown}

The fixed parameters in
\cref{eq:accelerated-momentum,eq:accelerated-gradient-step} use both $\beta$
and $\mu$. If certified bounds $B\ge\beta$ and $0<m\le\mu$ are available,
one may use $B$ and $m$ in their place; the resulting guarantee depends on
$B/m$. We now separate the two kinds of missing information.

First suppose that $\mu$ is known and $\beta$ is unknown. Section~\ref{sec:backtracking}
searched for a curvature estimate by testing the smooth quadratic upper bound
at a proposed gradient step. The same test can be applied at the extrapolated
point $y_t$. A trial curvature estimate $M>0$ defines
\begin{equation}
 \theta^+(M)=y_t-\frac1M\nabla L(y_t)
 \label{eq:accelerated-trial-step}
\end{equation}
and is accepted when
\begin{equation}
 L(\theta^+(M))
 \le L(y_t)+\langle\nabla L(y_t),\theta^+(M)-y_t\rangle
       +\frac{M}{2}\|\theta^+(M)-y_t\|_2^2.
 \label{eq:accelerated-model-test}
\end{equation}
The acceptance inequality is the smooth quadratic upper bound with $M$ in
place of $\beta$, evaluated only at $y_t$ and $\theta^+(M)$. Every
$M\ge\beta$ passes the test. After a rejection, replace $M$ by $\rho M$ for a
fixed $\rho>1$ and test the shorter step. The search therefore terminates
after finitely many objective evaluations.

Acceleration links the gradient step and the momentum coefficient in the same
potential inequality. An \emph{accelerated backtracking} method supplies the
accepted curvature estimate to both parts of the accelerated update. With
$\mu$ known and exact objective values available, standard accelerated
backtracking schemes attain
\begin{equation}
 N_\epsilon
 =O\left(\sqrt{\frac\beta\mu}
       \log\frac{L(\theta_0)-L^\star}{\epsilon}\right)
 \label{eq:accelerated-backtracking-rate}
\end{equation}
without receiving $\beta$ as input. Changing the gradient-step denominator
while keeping $\gamma$ fixed leaves the potential calculation in
\cref{eq:accelerated-design-contraction} unjustified; accelerated
backtracking recomputes the coupling coefficients when the curvature estimate
changes.

Unknown $\mu$ requires a different procedure because the momentum coefficient
cannot be obtained from the smoothness test. Suppose now that $L$ is convex,
its minimum is attained, and it is $\beta$-smooth and $\mu$-PL for an unknown
$\mu>0$. Write $\Theta^\star=\argmin L$ and define
\begin{equation}
 \operatorname{dist}(z,\Theta^\star)
 =\inf_{\theta^\star\in\Theta^\star}\|z-\theta^\star\|_2.
 \label{eq:distance-to-minimizers}
\end{equation}
The PL inequality implies the quadratic-growth bound
\begin{equation}
 L(z)-L^\star
 \ge\frac\mu2\operatorname{dist}(z,\Theta^\star)^2.
 \label{eq:pl-quadratic-growth}
\end{equation}
Run the smooth-convex accelerated method for $s$ iterations and then set both
stored parameter iterates equal to the current parameter before continuing.
The inner method can use a known $\beta$ or the accelerated backtracking
procedure above.
One such run between resets is a \emph{restart epoch}. If one epoch uses $s$
gradients,
\cref{eq:accelerated-restart-one-epoch} gives
\begin{equation}
 L(\mathsf A_s(z))-L^\star
 \le\frac{4\kappa}{(s+1)^2}\bigl(L(z)-L^\star\bigr).
 \label{eq:accelerated-restart-main}
\end{equation}
Thus an epoch length $s=\Theta(\sqrt\kappa)$ reduces the objective gap by a
constant factor, and repeated restarts attain
\begin{equation}
 N_\epsilon
 =O\left(\sqrt\kappa
 \log\frac{L(\theta_0)-L^\star}{\epsilon}\right).
 \label{eq:accelerated-restart-main-complexity}
\end{equation}
When $\kappa$ is unknown, an adaptive restart rule searches for an epoch
length on the scale $\sqrt\kappa$. For example, if $L^\star$ is available,
one can try epoch lengths $1,2,4,\ldots$ from the same starting point until
the observed objective gap decreases by one half. The successful search uses
$O(\sqrt\kappa)$ gradients, after which the accepted epoch length can be
reused. \Cref{ex:l12-accelerated-restart} asks you to prove this claim. Choosing
the epoch length without either $\mu$ or $L^\star$ requires a different test;
the bibliographical notes cite provable rules and common practical heuristics.
The optional reading derives the smooth-convex Lyapunov bound and the restart
calculation.

\section{Preconditioning changes the geometry}
\label{sec:preconditioning}

This section studies a different response to poor conditioning. Backtracking
chooses one scalar step, whereas preconditioning changes the relative scale
and direction of the parameter coordinates. A \emph{preconditioned gradient
update} uses a symmetric positive-definite matrix $P$:
\begin{equation}
 \theta_{t+1}=\theta_t-\eta P\nabla L(\theta_t).
 \label{eq:preconditioned-update}
\end{equation}
The matrix can rescale coordinates and mix them. The purpose of $P$ is to make
the objective better conditioned in the geometry used by the update.

Matrix preconditioning introduces a second meaning of adaptation. A line
search changes one scalar according to the objective decrease observed at the
current point. A
preconditioner changes the relative scale and direction of the coordinates.
The difference between scalar line search and matrix preconditioning matters
in a large model: one scalar can control overall stability, while a matrix can
assign substantially different effective stepsizes to different parameter
directions.

\begin{proposition}[Preconditioning as a change of coordinates]
\label{prop:preconditioning-coordinates}
Let $\theta=P^{1/2}z$ and define
$\widetilde L(z)=L(P^{1/2}z)$. If
\begin{equation}
 z_{t+1}=z_t-\eta\nabla\widetilde L(z_t),
 \label{eq:preconditioned-z-update}
\end{equation}
then $\theta_t=P^{1/2}z_t$ satisfies
\begin{equation}
 \theta_{t+1}=\theta_t-\eta P\nabla L(\theta_t).
 \label{eq:preconditioning-coordinate-equivalence}
\end{equation}
If $L$ is quadratic with Hessian $H$, then
the Hessian in the $z$ coordinates is
\begin{equation}
 \widetilde H=P^{1/2}HP^{1/2}.
 \label{eq:preconditioned-hessian}
\end{equation}
\end{proposition}

\begin{proof}
The chain rule gives
$\nabla\widetilde L(z)=P^{1/2}\nabla L(P^{1/2}z)$. Thus
\[
 z_{t+1}=z_t-\eta\nabla\widetilde L(z_t)
\]
maps, after multiplication by $P^{1/2}$, to
\[
 \theta_{t+1}=\theta_t-\eta P\nabla L(\theta_t).
\]
Applying the chain rule twice gives \cref{eq:preconditioned-hessian}.
\end{proof}

For the quadratic in \cref{eq:quadratic-objective}, the ideal choice
$P=H^{-1}$ gives $\widetilde H=I$ and condition number one. A dense global
matrix, however, requires $\Theta(d^2)$ storage and arithmetic per
matrix--vector product for a model with $d$ parameters. Such a cost is not a
realistic per-iteration option at modern model dimensions. A whole-model
preconditioner must normally cost $O(d)$, or at most about $O(d\log d)$, to
store and apply. More expensive factorizations can be useful only on small
blocks or when the factorization cost is amortized over many gradient steps.

Restricting the cost also restricts the matrices that can be represented. A diagonal
preconditioner costs $O(d)$ but cannot represent any cross-coordinate
interaction. A block-diagonal preconditioner with blocks of size at most $b$
costs $O(db)$ to store and apply; the block-diagonal form discards
interactions between blocks.
Factoring all such blocks can cost $O(db^2)$ when the factors are refreshed.
Low-rank and Kronecker-factored forms retain other selected interactions, but
neither type can represent an arbitrary dense inverse Hessian. If the
discarded interactions contain the poorly conditioned directions, no general
condition-number improvement follows. \Cref{ex:l12-rotated-preconditioner}
gives a two-dimensional example in which diagonal Hessian scaling leaves the
condition number unchanged while full preconditioning reduces it to one.

A common neural-network example makes the approximation explicit. The Fisher
information matrix is the expected outer product of log-likelihood gradients.
For a fully connected layer with weight matrix $W$, a per-example gradient has
the outer-product form
\[
 \nabla_W\ell=\delta a^\top,
\]
where $a$ is the incoming activation and $\delta$ is the back-propagated
derivative at the layer output. The corresponding Fisher block contains
\begin{equation}
\E\bigl[(aa^\top)\otimes(\delta\delta^\top)\bigr].
 \label{eq:exact-fisher-block}
\end{equation}
Here $\otimes$ denotes the Kronecker product.
A Kronecker-factored approximation replaces the expectation in
\cref{eq:exact-fisher-block} by
\begin{equation}
 \E[aa^\top]\otimes\E[\delta\delta^\top].
 \label{eq:kfac-block}
\end{equation}
The factorization discards dependence between activations and
back-propagated derivatives, and a layerwise block approximation also
discards cross-layer interactions. Rigorous convergence results exist for
restricted overparameterized or infinite-width network models, but no
general theorem says that a low-rank or Kronecker-factored preconditioner
improves the conditioning of a finite multilayer network.

Preconditioning also changes solution selection when the minimizer is not
unique. Lecture~10 showed that the optimization geometry can determine which
interpolating or separating solution is approached. Better conditioning and
changed implicit bias are therefore separate consequences of the same matrix
update.

\section{AdaGrad adapts to observed gradient geometry}
\label{sec:adagrad}

Backtracking adapts to curvature by evaluating the objective at trial points.
That test is practical when reliable objective values are inexpensive. In
large-scale minibatch training, evaluating the full empirical objective
requires a pass over the data, and every trial on the current minibatch
requires another forward pass. Trials evaluated on different minibatches also
produce noisy values for different functions. These costs motivate adaptive
rules that use the first-order observations already computed by the training
procedure.

AdaGrad gives a clean example of such a rule. Its guarantee applies to
subgradients of convex losses, including losses with no smoothness constant
$\beta$. This section first studies scalar adaptation in online convex
optimization: the method nearly matches the bound obtained by choosing the
best constant stepsize after all gradients have been observed, without
knowing the horizon or the future gradient norms. AdaGrad then applies the
same idea coordinate by coordinate. The coordinatewise refinement can improve
the scalar bound when gradient activity is sparse or uneven.

\paragraph{Scope of the online model.}
The online model supplies the historical and mathematically clean setting for
AdaGrad's adaptive subgradient guarantee. The regret theorem gives an exact
comparison with a stepsize selected after the gradient sequence is known.
Its formal scope is online convex optimization. Under an i.i.d.\ sampling
model with $L(x)=\E[\ell_t(x)]$, convexity and the averaged decision
$\overline x_T=T^{-1}\sum_{t=1}^T x_t$ give the online-to-batch relation
\begin{equation}
 \E[L(\overline x_T)-L(u)]
 \le\frac1T\E[\operatorname{Reg}_T(u)].
 \label{eq:online-to-batch}
\end{equation}
Language-model training is a nonconvex stochastic problem outside that scope.
Here the theorem identifies the origin and precise meaning of AdaGrad's
adaptive update.

The online problem supplies a sequence of convex losses instead of the one
fixed objective $L$ used above. To mark the change of setting, we write $x_t$
for the online decisions. At round $t$, the learner chooses $x_t$ before
seeing the current loss $\ell_t$ and then receives a subgradient
$g_t\in\partial\ell_t(x_t)$. Performance is compared with one fixed decision
$u$ chosen after the whole sequence is known. The regret is
\begin{equation}
 \operatorname{Reg}_T(u)
 =\sum_{t=1}^T\bigl(\ell_t(x_t)-\ell_t(u)\bigr).
 \label{eq:online-regret}
\end{equation}
A sublinear regret bound means that, for every fixed comparator $u$,
\begin{equation}
 \operatorname{Reg}_T(u)=o(T)
 \quad\Longleftrightarrow\quad
 \frac1T\sum_{t=1}^T
 \bigl(\ell_t(x_t)-\ell_t(u)\bigr)\longrightarrow0.
 \label{eq:sublinear-regret-meaning}
\end{equation}
The guarantee is deterministic and allows the loss sequence to be
adversarial; no sampling model is required.

Let $\cX$ be a closed convex set with Euclidean diameter at most $D>0$. Projected
online gradient descent with a constant stepsize $\eta>0$ satisfies
\begin{equation}
 \operatorname{Reg}_T(u)
 \le \frac{D^2}{2\eta}
      +\frac\eta2\sum_{t=1}^T\|g_t\|_2^2.
 \label{eq:fixed-eta-oco}
\end{equation}
If the entire sequence were known before the run and at least one gradient
were nonzero, the right side would be minimized by
\begin{equation}
 \eta^\star
 =\frac{D}{\sqrt{\sum_{t=1}^T\|g_t\|_2^2}}.
 \label{eq:hindsight-stepsize}
\end{equation}
The best value of the constant-stepsize upper bound is
\begin{equation}
 \inf_{\eta>0}
 \left\{
  \frac{D^2}{2\eta}
  +\frac\eta2\sum_{t=1}^T\|g_t\|_2^2
 \right\}
 =D\left(\sum_{t=1}^T\|g_t\|_2^2\right)^{1/2}.
 \label{eq:hindsight-bound-value}
\end{equation}
The tuned stepsize depends on future gradients and therefore cannot be
supplied to the online algorithm.

\begin{theorem}[Scalar adaptation to the observed gradient scale]
\label{thm:scalar-adaptation}
Define
\[
 S_t=\left(\sum_{s=1}^t\|g_s\|_2^2\right)^{1/2}
\]
and use the projected update
\begin{equation}
 x_{t+1}=\Pi_{\cX}\left(x_t-\frac{D}{S_t}g_t\right),
 \label{eq:scalar-adaptive-update}
\end{equation}
with no change when $S_t=0$. Then, for every $u\in\cX$,
\begin{equation}
 \operatorname{Reg}_T(u)
 \le\frac32D
      \left(\sum_{t=1}^T\|g_t\|_2^2\right)^{1/2}.
 \label{eq:scalar-adaptive-regret}
\end{equation}
\end{theorem}

\begin{proof}
If $S_T=0$, then every $g_t=0$, and convexity gives
$\ell_t(x_t)-\ell_t(u)\le0$ on every round. Suppose therefore that $S_T>0$,
and let $t_0$ be the first index for which $S_{t_0}>0$. Rounds before $t_0$
again contribute nonpositive regret. For $t\ge t_0$, set
$\eta_t=D/S_t$. Since $u\in\cX$, the projection inequality gives
\begin{equation}
 \|x_{t+1}-u\|_2^2
 \le \|x_t-\eta_tg_t-u\|_2^2
 =\|x_t-u\|_2^2
   -2\eta_t\langle g_t,x_t-u\rangle
   +\eta_t^2\|g_t\|_2^2.
 \label{eq:scalar-adaptive-projection-step}
\end{equation}
Rearranging \cref{eq:scalar-adaptive-projection-step} yields
\begin{equation}
 \langle g_t,x_t-u\rangle
 \le
 \frac{\|x_t-u\|_2^2-\|x_{t+1}-u\|_2^2}{2\eta_t}
 +\frac{\eta_t}{2}\|g_t\|_2^2.
 \label{eq:scalar-adaptive-inner-product-step}
\end{equation}
Convexity gives
\[
 \ell_t(x_t)-\ell_t(u)\le\langle g_t,x_t-u\rangle.
\]
Combining the last two inequalities gives the one-step bound
\[
 \ell_t(x_t)-\ell_t(u)
 \le \frac{\|x_t-u\|_2^2-\|x_{t+1}-u\|_2^2}{2\eta_t}
      +\frac{\eta_t}{2}\|g_t\|_2^2.
\]

We next sum the distance terms. Summation by parts is the discrete analogue
of integration by parts. The mnemonic
\[
 \int Fg=FG-\int fG,
 \qquad F'=f,\quad G'=g,
\]
says to move the derivative---or, for a sum, the finite difference---from one
factor to the other and retain the boundary terms. Put
$A_t=\|x_t-u\|_2^2$. Directly expanding the finite sum gives
\begin{align}
 \sum_{t=t_0}^T\frac{A_t-A_{t+1}}{2\eta_t}
 &=\frac{A_{t_0}}{2\eta_{t_0}}
   +\sum_{t=t_0+1}^T
    \left(\frac1{2\eta_t}-\frac1{2\eta_{t-1}}\right)A_t
   -\frac{A_{T+1}}{2\eta_T}.                       \label{eq:scalar-sbp}
\end{align}
The sequence $S_t$ is nondecreasing, so $1/\eta_t=S_t/D$ is also
nondecreasing. Every coefficient in the middle sum of
\cref{eq:scalar-sbp} is therefore nonnegative. Using $A_t\le D^2$ and
discarding the final nonpositive term yields
\begin{equation}
 \sum_{t=t_0}^T\frac{A_t-A_{t+1}}{2\eta_t}
 \le D^2\left[
       \frac1{2\eta_{t_0}}
       +\sum_{t=t_0+1}^T
        \left(\frac1{2\eta_t}-\frac1{2\eta_{t-1}}\right)
      \right]
 =\frac{D^2}{2\eta_T}
  =\frac{DS_T}{2}.                                  \label{eq:scalar-distance-sum}
\end{equation}

For the gradient terms, $S_t^2-S_{t-1}^2=\|g_t\|_2^2$ and hence
\begin{equation}
 \sum_{t=t_0}^T\frac{\|g_t\|_2^2}{S_t}
 =\sum_{t=t_0}^T
   (S_t-S_{t-1})\frac{S_t+S_{t-1}}{S_t}
 \le2\sum_{t=t_0}^T(S_t-S_{t-1})
  =2S_T.                                             \label{eq:scalar-gradient-sum}
\end{equation}
Multiplying \cref{eq:scalar-gradient-sum} by $D/2$ bounds the gradient
sum in the one-step inequality by $DS_T$. Combining the two sums proves
\cref{eq:scalar-adaptive-regret}.
\end{proof}

Combining \cref{eq:scalar-adaptive-regret,eq:hindsight-bound-value} gives
\begin{equation}
 \operatorname{Reg}_T(u)
 \le\frac32
 \inf_{\eta>0}
 \left\{
  \frac{D^2}{2\eta}
  +\frac\eta2\sum_{t=1}^T\|g_t\|_2^2
 \right\}.
 \label{eq:scalar-adaptation-versus-hindsight}
\end{equation}
Thus the adaptive guarantee is within a factor $3/2$ of the best value of the
constant-stepsize upper bound \cref{eq:fixed-eta-oco}. The procedure does not
need the horizon or a bound on the gradient norms, but the procedure still
uses the domain diameter $D$. Adaptation is always relative to a specified set
of unknown and known quantities.

Diagonal AdaGrad applies the same argument separately to each coordinate.
Let $\cX=\prod_{i=1}^d[a_i,b_i]$, with coordinate diameter
$D_i=b_i-a_i$.

Define the accumulated coordinate magnitude
\begin{equation}
 G_{t,i}=\left(\delta^2+\sum_{s=1}^t g_{s,i}^2\right)^{1/2},
 \qquad \delta\ge0.
 \label{eq:adagrad-accumulator}
\end{equation}
Diagonal AdaGrad uses
\begin{equation}
 x_{t+1,i}
 =\Pi_{[a_i,b_i]}
  \left(x_{t,i}-\frac{D_i}{G_{t,i}}g_{t,i}\right).
 \label{eq:adagrad-update}
\end{equation}
When $\delta=0$ and $G_{t,i}=0$, the coordinate is left unchanged. A small
positive $\delta$ is commonly used to avoid division by zero.

\begin{theorem}[Data-dependent regret of diagonal AdaGrad]
\label{thm:adagrad-regret}
With $\delta=0$, the iterates in \cref{eq:adagrad-update} satisfy, for every
$u\in\cX$,
\begin{equation}
 \sum_{t=1}^T\bigl(\ell_t(x_t)-\ell_t(u)\bigr)
 \le\frac32\sum_{i=1}^d D_i
       \left(\sum_{t=1}^Tg_{t,i}^2\right)^{1/2}.
 \label{eq:adagrad-regret}
\end{equation}
\end{theorem}

\Cref{ex:l12-adagrad-regret} asks you to prove the theorem. To compare the
diagonal and scalar guarantees, define
\[
 D_2=\left(\sum_{i=1}^dD_i^2\right)^{1/2},
 \qquad
 S_i=\left(\sum_{t=1}^Tg_{t,i}^2\right)^{1/2}.
\]
Cauchy--Schwarz gives
\begin{equation}
 \sum_{i=1}^dD_iS_i
 \le D_2\left(\sum_{t=1}^T\|g_t\|_2^2\right)^{1/2}.
 \label{eq:diagonal-versus-scalar}
\end{equation}
Thus the diagonal AdaGrad bound is never worse than the corresponding scalar
adaptive bound with the box's Euclidean diameter. The comparison is between
the two upper bounds; a comparison of realized losses requires separate
analysis. If all
$D_i$ are equal and only $k$ coordinates ever receive gradients of comparable
cumulative magnitude, then, in the exactly equal case,
\begin{equation}
 \frac{\sum_{i=1}^dD_iS_i}
 {D_2\bigl(\sum_{t=1}^T\|g_t\|_2^2\bigr)^{1/2}}
 =\sqrt{\frac{k}{d}}.
 \label{eq:adagrad-sparse-gain}
\end{equation}
The diagonal regret upper bound is therefore smaller by the factor
$\sqrt{d/k}$, with the largest factor $\sqrt d$ attained when $k=1$.

The theorem identifies the specific sense in which AdaGrad is adaptive.
AdaGrad records the gradient magnitudes observed in each coordinate, and the
regret bound depends on the same recorded quantities. A sparse or uneven
gradient sequence can therefore yield a smaller comparator bound than scalar
adaptation.

The smooth-objective and online analyses use gradient information
differently. For one fixed smooth objective, the quadratic upper bound
converts a gradient step into objective descent, giving an $O(1/T)$ convex
objective-gap bound or geometric objective-gap decay under the PL inequality.
In online convex optimization, adversarially changing losses prevent
telescoping of objective values from different rounds. The online proof uses a
squared-distance recursion and retains the cumulative squared gradients.
Under bounded gradients, the
theorem gives the usual average-regret bound: if $\|g_t\|_2\le G$, then
\begin{equation}
 \frac1T\sum_{t=1}^T\bigl(\ell_t(x_t)-\ell_t(u)\bigr)
 \le\frac{3DG}{2\sqrt T}.
 \label{eq:scalar-average-regret}
\end{equation}

AdaGrad pays for long-term memory in another way. The denominator
$G_{t,i}$ never decreases. A large gradient early in training permanently
reduces later steps in that coordinate. Permanent dependence on early
gradients is appropriate for the adversarial online regret theorem but can be
too conservative when gradient scales change during nonconvex training.

The statistic $G_{t,i}^2$ is a sum of squared gradients. A diagonal entry of
the Hessian is a derivative of a gradient coordinate. Particular
probabilistic models can relate these quantities, but their definitions and
general roles are different. AdaGrad constructs an adaptive diagonal geometry
from the observed gradients.

\section{Adapting to changing gradient scales}
\label{sec:adam}

AdaGrad retains every past squared gradient in its denominator. A large
gradient therefore reduces every later step in the same coordinate, even
after the relevant gradient scale has changed. Neural-network gradients can
also differ greatly across coordinates and across stages of training. Adam
uses recent squared gradients to set coordinatewise scales and uses recent
gradients to set the update direction. This section first explains that
design and then studies an exact logistic-regression example.

Given a stochastic gradient $g_t\in\R^d$ at $\theta_t$ and initial states
$m_0=v_0=0$, Adam uses
\begin{align}
 m_t&=\beta_1m_{t-1}+(1-\beta_1)g_t,
 \label{eq:adam-first-moment}\\
 v_t&=\beta_2v_{t-1}+(1-\beta_2)(g_t\odot g_t),
 \label{eq:adam-second-moment}\\
 \widehat m_t&=\frac{m_t}{1-\beta_1^t},
 \qquad
 \widehat v_t=\frac{v_t}{1-\beta_2^t},
 \label{eq:adam-bias-correction}\\
 \theta_{t+1}&=\theta_t-\alpha
 \frac{\widehat m_t}{\sqrt{\widehat v_t}+\varepsilon}.
 \label{eq:adam-update}
\end{align}
Products, square roots, and division are coordinatewise. The numerator
$\widehat m_t$ averages recent gradients and supplies the direction. The
denominator supplies one scale per coordinate. Unrolling its recursion gives
\begin{equation}
 v_t=(1-\beta_2)\sum_{s=1}^t\beta_2^{t-s}(g_s\odot g_s).
 \label{eq:adam-exponential-memory}
\end{equation}
Thus a squared gradient observed $t-s$ iterations ago receives weight
$\beta_2^{t-s}$. AdaGrad assigns weight one to every past squared gradient;
Adam assigns geometrically decreasing weights. The parameters
$\beta_1,\beta_2\in[0,1)$ determine the two memory timescales,
$\alpha>0$ is a base stepsize, and $\varepsilon\ge0$ controls the update when
the recent squared gradients are small. When $\varepsilon=0$, the coordinate
$i$ update in \cref{eq:adam-update} is defined whenever
$\widehat v_{t,i}>0$; the scalar examples below have a nonzero gradient at
every finite iterate.

The factors $1-\beta_1^t$ and $1-\beta_2^t$ correct the initialization at
zero. For a constant gradient sequence, the corrected quantities equal the
corresponding first and squared gradients at every iteration. The quantities
$\alpha$, $\beta_1$, $\beta_2$, and $\varepsilon$ remain algorithmic choices.

RMSProp is the earlier method that uses the same recent-square scaling with
the current gradient in the numerator. A common version uses
\begin{equation}
 \theta_{t+1}=\theta_t-\alpha
 \frac{g_t}{\sqrt{v_t}+\varepsilon},
 \qquad
 v_t=\beta_2v_{t-1}+(1-\beta_2)(g_t\odot g_t).
 \label{eq:rmsprop-update}
\end{equation}
Setting $\beta_1=0$ in Adam gives the bias-corrected version of this update.
Convergence-to-stationarity guarantees for RMSProp require explicit
smoothness, noise, and hyperparameter conditions. Convex counterexamples rule
out an unconditional convergence guarantee.

We next examine Adam's behavior on a simple separable logistic problem.
For the special case $\beta_1=\beta_2=0$, the calculation gives an exact
comparison between the normalized update and constant-stepsize gradient
descent, and shows that normalization can produce faster loss decay in this
setting.
\begin{example}[Loss decay on a one-dimensional logistic problem]
\label{ex:logistic-normalization}
Let
\begin{equation}
 L(\theta)=\log(1+e^{-\theta}),
 \qquad
 L'(\theta)=-\frac1{1+e^\theta}.
 \label{eq:scalar-logistic-loss}
\end{equation}
Constant-stepsize gradient descent satisfies
\begin{equation}
 \theta_{t+1}=\theta_t+\frac\eta{1+e^{\theta_t}}.
 \label{eq:scalar-logistic-gd}
\end{equation}
For every $\eta>0$, the iterates obey
\begin{equation}
 \lim_{t\to\infty}tL(\theta_t)=\frac1\eta.
 \label{eq:scalar-logistic-gd-rate}
\end{equation}
The normalized update with $\alpha>0$ gives
\begin{equation}
 \theta_{t+1}
 =\theta_t-\alpha\frac{L'(\theta_t)}{|L'(\theta_t)|}
  =\theta_t+\alpha,
 \qquad
 \lim_{t\to\infty}e^{\alpha t}L(\theta_t)
 =e^{-\theta_0}.
 \label{eq:scalar-logistic-normalized-rate}
\end{equation}
Thus the loss decays as $1/t$ for the constant gradient-descent stepsize and
exponentially for the normalized update.
\end{example}

To verify \cref{eq:scalar-logistic-gd-rate}, set
$u_t=e^{\theta_t}$. The iterates in \cref{eq:scalar-logistic-gd} increase
without bound, and
\[
 u_{t+1}-u_t
 =u_t\left[
   \exp\left(\frac\eta{1+u_t}\right)-1
  \right]\longrightarrow\eta.
\]
The Stolz--Ces\`aro theorem gives $u_t/t\to\eta$, and
$\log(1+1/u_t)\sim1/u_t$ proves
\cref{eq:scalar-logistic-gd-rate}. The identity
$\theta_t=\theta_0+\alpha t$ proves
\cref{eq:scalar-logistic-normalized-rate}.

The normalized update in the example is the RMSProp and Adam special case
$\beta_1=\beta_2=0$ and $\varepsilon=0$. The calculation proves exponential
loss decay for this scalar problem. Comparisons on other objectives require
their own analysis.

The same special case gives an exact interpretation of
$\varepsilon>0$. The update becomes
\begin{equation}
 \theta_{t+1}
 =\theta_t+\alpha
  \frac{|L'(\theta_t)|}{|L'(\theta_t)|+\varepsilon}.
 \label{eq:scalar-logistic-epsilon-update}
\end{equation}
For $0<\varepsilon<1$, the gradient threshold
$|L'(\theta)|=\varepsilon$ corresponds exactly to the loss level
\begin{equation}
 L(\theta)=-\log(1-\varepsilon).
 \label{eq:scalar-logistic-epsilon-level}
\end{equation}
When $|L'(\theta_t)|\gg\varepsilon$, the update in
\cref{eq:scalar-logistic-epsilon-update} is approximately the normalized
update. When $|L'(\theta_t)|\ll\varepsilon$, the update is approximately a
gradient step with stepsize $\alpha/\varepsilon$. In fact,
\begin{equation}
 \lim_{t\to\infty}tL(\theta_t)=\frac{\varepsilon}{\alpha}.
 \label{eq:scalar-logistic-epsilon-rate}
\end{equation}
To obtain the limit, again set $u_t=e^{\theta_t}$ and observe that
$u_{t+1}-u_t\to\alpha/\varepsilon$. Thus $\varepsilon=0$ preserves
exponential decay, whereas every fixed $\varepsilon>0$ eventually returns to
$1/t$ decay after a potentially long normalized phase. For general Adam, the
transition occurs when $\sqrt{\widehat v_{t,i}}$ is comparable to
$\varepsilon$. Each coordinate has its own transition, determined by the
recent gradient history.

Exponential forgetting removes the monotonicity used in AdaGrad's regret
proof. An old large gradient can cease to constrain the effective step even
when that gradient remains important for the optimization objective.

\begin{example}[An extreme forgetting failure]
\label{ex:adam-forgetting}
On $[-1,1]$, repeat the three linear losses
\[
 \ell_{3k+1}(x)=Cx,
 \qquad
 \ell_{3k+2}(x)=\ell_{3k+3}(x)=-x,
 \qquad C>2.
\]
The three losses sum to $(C-2)x$ over each block, so the best fixed point is
$x=-1$. Take the projected Adam special case $\beta_1=\beta_2=0$ and
$\varepsilon=0$, and fix $0<\alpha\le2$. For
$g_t=\ell_t'(x_t)$, the update is
\begin{equation}
 x_{t+1}=\Pi_{[-1,1]}
 \left(x_t-\alpha\frac{g_t}{|g_t|}\right).
 \label{eq:adam-forgetting-update}
\end{equation}
The update normalizes every nonzero gradient to its sign. Starting from
$x=1$, the iterates used for one block are
$1,1-\alpha,1$. The next block again begins at $1$. Every block therefore
incurs the same positive regret. After $K$ blocks,
\begin{equation}
 \sum_{t=1}^{3K}\bigl(\ell_t(x_t)-\ell_t(-1)\bigr)
 =K\bigl(2(C-2)+\alpha\bigr).
 \label{eq:adam-linear-regret}
\end{equation}
\Cref{ex:l12-adam-forgetting} asks you to calculate the iterates and the exact
block regret.
\end{example}

The example proves an online convex optimization failure:
\begin{equation}
 \frac{\operatorname{Reg}_{3K}(-1)}{3K}
 =\frac{2(C-2)+\alpha}{3}>0.
 \label{eq:adam-nonvanishing-average-regret}
\end{equation}
Failures also occur in stochastic convex optimization and in deterministic
batch optimization. Separate counterexamples establish the following two
statements.
\begin{enumerate}[(a),leftmargin=*]
\item Let $\Xi$ be a random data example. There is a convex loss
$\ell(\theta;\Xi)$ for which Adam applied to the stochastic objective
\begin{equation}
 L(\theta)=\E[\ell(\theta;\Xi)]
 \label{eq:adam-stochastic-objective}
\end{equation}
produces iterates satisfying
\begin{equation}
 \operatorname{dist}\bigl(\theta_t,\argmin_\theta L(\theta)\bigr)
 \not\longrightarrow0.
 \label{eq:adam-stochastic-nonconvergence}
\end{equation}
\item A deterministic batch version of Adam with bias correction has
trajectories that approach a nonoptimal period-two limit cycle on the strongly
convex quadratic
\begin{equation}
 L(\theta)=\frac c2\theta^2,
 \qquad c>0.
 \label{eq:adam-quadratic-objective}
\end{equation}
For such a trajectory, there are two distinct states
$s_+=(a,m_+,v_+)$ and $s_-=(-a,m_-,v_-)$ with $a\ne0$ such that
\begin{equation}
 (\theta_{2t},m_{2t},v_{2t})\longrightarrow s_+,
 \qquad
 (\theta_{2t+1},m_{2t+1},v_{2t+1})\longrightarrow s_-.
 \label{eq:adam-two-cycle}
\end{equation}
\end{enumerate}
The deterministic construction uses $\sqrt{v_t+\varepsilon}$ in the
denominator; the update in \cref{eq:adam-update} uses
$\sqrt{v_t}+\varepsilon$. The construction rules out a global convergence
guarantee for the former batch Adam variant. Its conclusion concerns the
existence of particular trajectories and leaves behavior from standard
initializations and default hyperparameters open. Long-memory variants such
as AMSGrad restore a nondecreasing second-moment envelope for the corresponding
convex analysis.

Positive convergence theorems for Adam provide basic reassurance under
restricted assumptions. An advantage over scalar adaptation or SGD requires
a sharper comparison. Recent results identify settings in which the
coordinatewise normalization can improve the bound. Examples
include objectives with nonuniform smoothness of the schematic form
\[
 \|\nabla L(u)-\nabla L(v)\|_2
 \lesssim
 \bigl(L_0+L_1\|\nabla L(u)\|_2\bigr)
 \|u-v\|_2
\]
locally, and problems whose gradients or curvature are favorable in
$\ell_\infty$ or blockwise geometry. These results give problem-dependent
comparisons with scalar methods. A universal ordering of Adam and tuned SGD
and a hindsight-best diagonal guarantee comparable to AdaGrad's theorem
remain unavailable.

Adam remains a practically important optimizer, including for transformer
training. Existing theory and empirical evidence support particular
mechanisms and regimes. The denominator is an exponential average of squared
gradients. Default hyperparameters still require empirical validation, and
the available theorems apply to specified objective and noise classes.

\subsection{Decoupling weight decay from adaptive scaling}

The logistic example sends the parameter to infinity along the separating
direction. Practical training often includes weight decay to control parameter
growth. For an adaptive method, the placement of this shrinkage changes the
update. This subsection derives two common updates and then returns to
the scalar logistic example.

Suppose an adaptive method uses a positive diagonal matrix $D_t$ and applies
an $\ell_2$ penalty to the objective:
\[
 F(\theta)=L(\theta)+\frac\lambda2\|\theta\|_2^2.
\]
The update is
\begin{equation}
 \theta_{t+1}
 =\theta_t-\alpha D_t^{-1}(\nabla L(\theta_t)+\lambda\theta_t)
 =(I-\alpha\lambda D_t^{-1})\theta_t
   -\alpha D_t^{-1}\nabla L(\theta_t).
 \label{eq:adaptive-l2-update}
\end{equation}
The adaptive matrix scales the penalty gradient. Coordinates with different
entries of $D_t$ receive different amounts of shrinkage.

AdamW applies \emph{decoupled weight decay}:
\begin{equation}
 \theta_{t+1}
 =(1-\alpha\lambda)\theta_t
  -\alpha D_t^{-1}\nabla L(\theta_t).
 \label{eq:adamw-update}
\end{equation}
The factor multiplying $\theta_t$ is independent of the adaptive gradient
geometry. For ordinary unpreconditioned gradient descent, $D_t=I$ and the two
updates coincide. For a general adaptive diagonal matrix, the two updates
define different algorithms. \Cref{ex:l12-adamw-decoupling} asks you to verify
the distinction.

For the normalized scalar logistic update in
\cref{ex:logistic-normalization}, AdamW gives
\begin{equation}
 \theta_{t+1}=(1-\alpha\lambda)\theta_t+\alpha.
 \label{eq:scalar-logistic-adamw}
\end{equation}
If $0<\alpha\lambda<2$, then
\begin{equation}
 \theta_t\longrightarrow\frac1\lambda,
 \qquad
 L(\theta_t)\longrightarrow\log(1+e^{-1/\lambda})>0.
 \label{eq:scalar-logistic-adamw-limit}
\end{equation}
The limiting parameter is finite, and the limiting unregularized logistic
loss is positive. The weight-decayed and unregularized procedures therefore
optimize different training criteria.

The algebra also clarifies the design choice. An $\ell_2$ penalty belongs to
the objective, so an adaptive optimizer preconditions the penalty gradient
together with the loss gradient. Decoupled weight decay prescribes uniform
multiplicative shrinkage in the chosen parameterization. The two updates
encode different parameter preferences.

\subsection{How acceleration and adaptation enter neural-network training}

The methods in this lecture solve different problems, and practical training
often combines several of them. Momentum SGD and Nesterov-style momentum are
used for neural networks. The smooth strongly convex acceleration theorem
assumes a deterministic objective with global constants $\mu$ and $\beta$.
General neural-network training has a nonconvex objective and noisy minibatch
gradients, so the theorem supplies no rate for that setting. The momentum
coefficient and the base learning rate therefore remain empirical choices in
such applications.

AdaGrad and Adam modify a different part of the update. AdaGrad assigns
coordinatewise scales from all past squared gradients. The permanent
accumulator is useful when rare or sparse coordinates should receive larger
updates, but the same permanence can make later updates too small when the
gradient scale changes. Adam uses exponential memory so that the
coordinatewise scales can change. AdamW combines that normalization with
uniform weight decay. Transformer training provides extensive empirical use
of AdamW. Available theory gives guarantees in restricted models and leaves a
universal comparison with momentum SGD open.

A learning-rate schedule controls the global multiplier left outside these
direction and normalization choices. For a known training budget $T$, a
single cosine decay has the form
\begin{equation}
 \alpha_t
 =\alpha_{\min}
  +\frac{\alpha_{\max}-\alpha_{\min}}2
   \left(1+\cos\frac{\pi t}{T}\right),
 \qquad 0\le t\le T.
 \label{eq:cosine-learning-rate}
\end{equation}
The schedule makes relatively large updates early and reduces the global
step scale smoothly near the end of a fixed-budget run. A \emph{warmup phase}
increases the learning rate from a smaller initial value to
$\alpha_{\max}$ during a prescribed number of initial iterations and is often
placed before the decay. Cosine decay is fixed before training:
$\alpha_t$ depends only on $t$ and the chosen schedule parameters. The choices
of $\alpha_{\max}$,
$\alpha_{\min}$, the warmup length, and $T$ still require specification.
When cosine decay is used with Adam, the coefficients $\beta_1$ and $\beta_2$
usually remain fixed. The weights $\beta_1^{t-s}$ and $\beta_2^{t-s}$ assigned
to past gradients therefore decay at the same rate throughout training. Near
the end of training, however, the smaller $\alpha_t$ means that the parameters
move less during one fixed-length memory window, so the remembered gradients
come from more similar parameter values. Forgetting per iteration is constant;
the memory window becomes longer relative to the rate at which the parameters
move.
The original cosine-annealing proposal used repeated cycles with warm
restarts; a single decay uses one cycle without restarting.

\section{How much tuning can be removed?}
\label{sec:tuning-free}

Each method studied above removes a particular tuning requirement.
Backtracking selects a safe stepsize without being supplied with $\beta$;
scalar and diagonal AdaGrad select stepsizes without knowing the future
gradient scales. RMSProp and Adam track recent coordinatewise scales while
retaining a base learning rate and memory parameters.
\Cref{ex:l12-tuned-comparisons} asks you to recover the exact comparisons with
the corresponding methods tuned in advance. This section addresses the
broader question: can one algorithm match a tuned method when several
properties of the problem are unknown?

We use the phrase \emph{tuning-free} for a comparison between a self-tuning
method and a specified reference method. A tuning-free claim must specify the
quantities supplied to the reference method, the looser information supplied
to the self-tuning method, and the performance criterion. Two important
results show both the scope and the limits of such claims.

\subsection{Adapting to the regularity of a deterministic objective}

Universal convex methods address an entire family of regularity regimes.
Suppose that $L$ is convex and that, for some unknown $\nu\in[0,1]$ and
$H_\nu>0$, the first-order oracle obeys
\begin{equation}
 \|g(u)-g(v)\|_2
 \le H_\nu\|u-v\|_2^\nu,
 \qquad g(u)\in\partial L(u).
 \label{eq:holder-first-order-oracle}
\end{equation}
Let $R=\|\theta_0-\theta^\star\|_2$. Given a target objective error
$\epsilon>0$, a universal accelerated method can attain
\begin{equation}
 L(\theta_N)-L^\star\le\epsilon
 \quad\text{after}\quad
 N=O\left(
   \inf_{\nu\in[0,1]}
   \left(\frac{H_\nu}{\epsilon}\right)^{\frac{2}{1+3\nu}}
   R^{\frac{2+2\nu}{1+3\nu}}
 \right)
 \label{eq:universal-convex-complexity}
\end{equation}
first-order oracle calls, without being supplied with $\nu$ or $H_\nu$.
For $\nu=1$, the bound becomes
$O(\sqrt{H_1R^2/\epsilon})$, the accelerated smooth-convex rate. For
$\nu=0$, it becomes $O(H_0^2R^2/\epsilon^2)$, the nonsmooth rate. The same
method selects the appropriate rate from the observed first-order behavior.

\subsection{Matching tuned SGD with several unknown scales}

Stochastic optimization makes the comparison more demanding because the
observations used for tuning are noisy. Consider a convex,
$\beta$-smooth objective on a convex set of diameter $D$. At a query
$\theta$, suppose the oracle returns $G(\theta)$ satisfying
\begin{equation}
 \E[G(\theta)\mid\theta]=\nabla L(\theta),
 \qquad
 \|G(\theta)-\nabla L(\theta)\|_2\le\sigma
 \quad\text{almost surely}.
 \label{eq:tuning-free-oracle}
\end{equation}
With suitable tuning, variants of SGD attain the high-probability rate
\begin{equation}
 L(\widehat\theta_T)-L^\star
 =\widetilde O\left(
    \frac{\beta D^2}{T}+\frac{\sigma D}{\sqrt T}
  \right),
 \label{eq:tuned-sgd-stochastic-rate}
\end{equation}
where $\widetilde O$ suppresses logarithmic factors in the failure probability
and the number of oracle calls. A tuning-free counterpart of SGD receives
loose ranges containing $\beta$, $D$, and $\sigma$, and must match
\cref{eq:tuned-sgd-stochastic-rate} up to logarithmic factors in the widths of
those ranges.

\begin{proposition}[Scope of tuning-free stochastic convex optimization]
\label{prop:tuning-free-scope}
For the bounded-domain problem above, algorithms exist that match
\cref{eq:tuned-sgd-stochastic-rate} up to logarithmic factors without the exact
values of $\beta$, $D$, and $\sigma$. On the unbounded domain $\R^d$, replace
$D$ by the initial distance
\begin{equation}
 D_\star=\|\theta_0-\theta^\star\|_2.
 \label{eq:initial-distance-tuning}
\end{equation}
No algorithm supplied only with loose ranges for the problem scales can
guarantee
\begin{equation}
 L(\widehat\theta_T)-L^\star
 =\widetilde O\left(
    \frac{\beta D_\star^2}{T}
    +\frac{\sigma D_\star}{\sqrt T}
  \right)
 \label{eq:unbounded-tuning-free-target}
\end{equation}
uniformly over all smooth convex objectives and bounded-noise stochastic
oracles. The impossibility already holds for one-dimensional objectives.
\end{proposition}

The logarithmic factors depend on the loose parameter ranges, $T$, and the
requested failure probability.

The bounded domain supplies a finite upper bound on how far a minimizer can be
from the initial point. On an unbounded domain, two stochastic
problems can produce the same observations with substantial probability while
having different values of $D_\star$ and different optimal tuning. An
algorithm then makes the same choices on both problems, although matching
\cref{eq:unbounded-tuning-free-target} would require different choices. Extra
assumptions on the noise can make its scale estimable and restore positive
results. The bibliographical notes give the precise theorems and references.

For neural-network training, the result identifies the information that a
strong tuning claim would have to address. The parameter domain is usually
unbounded, the objective is nonconvex, and the minibatch oracle is noisy.
Adam's coordinatewise normalization leaves the base learning-rate schedule,
memory parameters, numerical stabilizer, and weight decay to be selected.
A theorem comparing an automatically selected training recipe with the best
tuned AdamW or SGD recipe would require a specified model class, stochastic
oracle, training budget, and performance measure.

\section{Optional reading\suppmark}
\label{sec:optional}

This optional section proves the smooth-convex accelerated rate with a
time-varying Lyapunov function, uses that rate to obtain linear convergence by
restarting, and gives a spectral explanation for quadratic acceleration.

\subsection{A Lyapunov proof of smooth-convex acceleration}
\label{sec:convex-acceleration-lyapunov}

The smooth-convex accelerated method uses a growing coefficient in its
Lyapunov function. The growth of that coefficient produces the $1/t^2$ rate
that will be converted into geometric decay by restarting.

Let $A_0=0$ and initialize $\theta_0=v_0=z$. For the state
$(\theta_t,v_t)$, consider
\begin{equation}
 E_t=A_t\bigl(L(\theta_t)-L^\star\bigr)
 +\frac12\|v_t-\theta^\star\|_2^2.
 \label{eq:convex-acceleration-potential}
\end{equation}
The design objective is to keep $E_t$ nonincreasing while making $A_t$ grow.
For $v_{t+1}=v_t-a_{t+1}g_t$, expansion gives
\begin{equation}
 \frac12\|v_{t+1}-\theta^\star\|_2^2
 =\frac12\|v_t-\theta^\star\|_2^2
 -a_{t+1}\langle g_t,v_t-\theta^\star\rangle
 +\frac{a_{t+1}^2}{2}\|g_t\|_2^2.
 \label{eq:convex-acceleration-state-step}
\end{equation}
For $\theta_{t+1}=y_t-g_t/\beta$, smoothness gives
\begin{equation}
 A_{t+1}\bigl(L(\theta_{t+1})-L^\star\bigr)
 \le A_{t+1}\bigl(L(y_t)-L^\star\bigr)
 -\frac{A_{t+1}}{2\beta}\|g_t\|_2^2.
 \label{eq:convex-acceleration-smooth-step}
\end{equation}
The two squared-gradient terms cancel when
\begin{equation*}
 \frac{A_{t+1}}{2\beta}=\frac{a_{t+1}^2}{2}
 \qquad\Longleftrightarrow\qquad
 A_{t+1}=\beta a_{t+1}^2.
\end{equation*}
Use $a_{t+1}>0$ also as the increment $A_{t+1}-A_t$. The two design
requirements are therefore
\begin{equation}
 A_{t+1}=A_t+a_{t+1}=\beta a_{t+1}^2,
 \label{eq:convex-acceleration-weights}
\end{equation}
and define
\begin{align}
 y_t&=\frac{A_t\theta_t+a_{t+1}v_t}{A_{t+1}},\notag\\
 g_t&=\nabla L(y_t),\notag\\
 v_{t+1}&=v_t-a_{t+1}g_t,\notag\\
 \theta_{t+1}&=\frac{A_t\theta_t+a_{t+1}v_{t+1}}{A_{t+1}}
 =y_t-\frac1\beta g_t.
 \label{eq:convex-accelerated-coupling}
\end{align}

\begin{proposition}[Smooth-convex acceleration]
\label{prop:smooth-convex-acceleration}
Suppose that $L$ is convex and $\beta$-smooth and that the minimizer set
$\Theta^\star$ is nonempty. The iterates in
\cref{eq:convex-accelerated-coupling} satisfy $E_{t+1}\le E_t$. Consequently,
for every $t\ge1$,
\begin{equation}
 L(\theta_t)-L^\star
 \le\frac{2\beta\operatorname{dist}(z,\Theta^\star)^2}{(t+1)^2}.
 \label{eq:convex-acceleration-lyapunov-rate}
\end{equation}
\end{proposition}

\begin{proof}
The smoothness and auxiliary-state terms have already been expanded in
\cref{eq:convex-acceleration-smooth-step,eq:convex-acceleration-state-step}.
After substituting $A_{t+1}=\beta a_{t+1}^2$ into the smoothness bound, it
remains to control the term $L(y_t)-L^\star$.
Convexity at $y_t$, applied to $\theta_t$ and $\theta^\star$ with weights
$A_t$ and $a_{t+1}$, respectively, gives
\begin{align*}
 &A_t\bigl(L(\theta_t)-L(y_t)\bigr)
 +a_{t+1}\bigl(L^\star-L(y_t)\bigr)\\
 &\quad\ge\left\langle g_t,
 A_t(\theta_t-y_t)+a_{t+1}(\theta^\star-y_t)\right\rangle
 =-a_{t+1}\langle g_t,v_t-\theta^\star\rangle,
\end{align*}
where the equality uses
$A_{t+1}y_t=A_t\theta_t+a_{t+1}v_t$. Rearranging yields
\begin{equation}
 A_{t+1}\bigl(L(y_t)-L^\star\bigr)
 \le A_t\bigl(L(\theta_t)-L^\star\bigr)
 +a_{t+1}\langle g_t,v_t-\theta^\star\rangle.
 \label{eq:convex-acceleration-convex-step}
\end{equation}
Adding
\cref{eq:convex-acceleration-smooth-step,eq:convex-acceleration-convex-step,eq:convex-acceleration-state-step}
cancels the inner-product and squared-gradient terms and proves
$E_{t+1}\le E_t$.

The rate now follows from a lower bound on $A_t$. If
$q_t=\sqrt{\beta A_t}$, then
\cref{eq:convex-acceleration-weights} implies
$q_{t+1}^2=q_t^2+q_{t+1}$, and hence
$q_{t+1}\ge q_t+1/2$. Since $q_1=1$,
\begin{equation}
 A_t\ge\frac{(t+1)^2}{4\beta}
 \qquad(t\ge1).
 \label{eq:convex-acceleration-weight-growth}
\end{equation}
Finally, choose a minimizer $\theta^\star$ closest to $z$. Then
$E_t\le E_0=\operatorname{dist}(z,\Theta^\star)^2/2$ and
$A_t(L(\theta_t)-L^\star)\le E_t$ give
\cref{eq:convex-acceleration-lyapunov-rate}.
\end{proof}

The cancellation in the proof explains the coefficient equation
$A_{t+1}=\beta a_{t+1}^2$. The method is a standard coupled representation of
Nesterov's smooth-convex accelerated gradient method; eliminating $v_t$ gives
its familiar two-step form.

\subsection{Restarting acceleration under convex PL}
\label{sec:restart-optional}

This subsection derives the accelerated linear rate for smooth convex PL
objectives. The argument combines smooth-convex acceleration with the
quadratic-growth consequence of PL. Strong convexity is one sufficient
condition for this result.

Let $\mathsf A_s(z)$ denote the output $\theta_s$ of
\cref{eq:convex-accelerated-coupling} after $s$ gradient evaluations, starting
from $z$. By \cref{prop:smooth-convex-acceleration},
\begin{equation}
 L(\mathsf A_s(z))-L^\star
 \le\frac{2\beta\operatorname{dist}(z,\Theta^\star)^2}{(s+1)^2}.
 \label{eq:accelerated-convex-rate}
\end{equation}
The quadratic-growth bound \cref{eq:pl-quadratic-growth} gives
\begin{equation}
 \operatorname{dist}(z,\Theta^\star)^2
 \le\frac2\mu\bigl(L(z)-L^\star\bigr).
 \label{eq:quadratic-growth-distance-gap}
\end{equation}
Combining the two inequalities yields
\begin{equation}
 L(\mathsf A_s(z))-L^\star
 \le\frac{4\kappa}{(s+1)^2}
       \bigl(L(z)-L^\star\bigr).
 \label{eq:accelerated-restart-one-epoch}
\end{equation}
Choose $s=\lceil\sqrt{8\kappa}\rceil-1$ and restart the method from the
preceding output after every $s$ gradients. If $z_r$ is the parameter after
$r$ epochs, then
\begin{equation}
 L(z_r)-L^\star
 \le 2^{-r}\bigl(L(z_0)-L^\star\bigr).
 \label{eq:accelerated-restart-linear-rate}
\end{equation}
The total number of gradient evaluations needed to obtain
$L(z_r)-L^\star\le\epsilon$ is therefore
\begin{equation}
 O\left(\sqrt\kappa
 \log\frac{L(z_0)-L^\star}{\epsilon}\right).
 \label{eq:accelerated-restart-complexity}
\end{equation}
The additional symbols $\mathsf A_s$, $s$, and $r$ describe the inner method,
the epoch length, and the epoch index. They are needed only for this restart
construction.

\subsection{Implementing the minimax polynomial by Chebyshev iteration}
\label{sec:chebyshev}

This subsection identifies the degree-$t$ polynomial with the smallest
worst-case residual on $[\mu,\beta]$ and derives a first-order iteration that
generates that polynomial. Chebyshev semi-iteration has the algebraic form of
a heavy-ball update with time-varying coefficients. It is neither the
fixed-coefficient Polyak heavy-ball method nor the Nesterov method in
\cref{sec:acceleration}.

For the quadratic \cref{eq:quadratic-objective}, a linear first-order method
whose iterates use past gradients and linear combinations has an error of the
form
\[
 e_t=p_t(H)e_0,
 \qquad p_t(0)=1,
\]
for a polynomial $p_t$ of degree at most $t$. The condition $p_t(0)=1$
has a direct meaning: when $\lambda=0$, every gradient vanishes in that
eigendirection, so the error remains equal to its initial value.
For every $H$ satisfying $\mu I\preceq H\preceq\beta I$, the Euclidean error
satisfies
\begin{equation}
 \|e_t\|_2
 \le
 \max_{\lambda\in[\mu,\beta]}|p_t(\lambda)|\,\|e_0\|_2.
 \label{eq:polynomial-worst-case}
\end{equation}

\begin{proposition}[Polynomials generated by gradient descent]
\label{prop:gd-polynomials}
Constant-step gradient descent generates
\[
 p_t(\lambda)=(1-\eta\lambda)^t.
\]
Gradient descent with predetermined steps $\eta_0,\ldots,\eta_{t-1}$
generates
\[
 p_t(\lambda)=\prod_{j=0}^{t-1}(1-\eta_j\lambda).
\]
Thus repeated constant-step gradient descent restricts $p_t$ to a power of
one linear factor.
\end{proposition}

\begin{proof}
Each eigencoordinate satisfies
$e_{j+1}(\lambda)=(1-\eta_j\lambda)e_j(\lambda)$. Iterating the scalar
recursion gives both formulas.
\end{proof}

Let $T_t$ be the degree-$t$ Chebyshev polynomial, characterized by
$T_t(\cos\phi)=\cos(t\phi)$.
Assume $0<\mu<\beta$ throughout this subsection. If $\mu=\beta$, then
$H=\beta I$, and one gradient step with stepsize $1/\beta$ reaches
$\theta^\star$.

\begin{theorem}[Minimax residual polynomial]
\label{thm:chebyshev-minimax}
For every integer $t\ge1$,
\begin{equation}
 \min_{\substack{p\in\R[\lambda],\ \deg(p)\le t\\p(0)=1}}
 \max_{\lambda\in[\mu,\beta]}|p(\lambda)|
 =
 \frac{1}{\left|T_t\left(\frac{\beta+\mu}{\beta-\mu}\right)\right|}.
 \label{eq:chebyshev-minimax-value}
\end{equation}
A minimizer is
\begin{equation}
 p_t^\star(\lambda)
 =\frac{T_t\left(\frac{\beta+\mu-2\lambda}{\beta-\mu}\right)}
        {T_t\left(\frac{\beta+\mu}{\beta-\mu}\right)}.
 \label{eq:chebyshev-polynomial}
\end{equation}
\end{theorem}

\begin{proof}
The affine map
\[
 x=\frac{\beta+\mu-2\lambda}{\beta-\mu}
\]
sends $[\mu,\beta]$ to $[-1,1]$ and sends $\lambda=0$ to
$a=(\beta+\mu)/(\beta-\mu)>1$. Under the displayed affine change of variable,
a feasible polynomial $p$ becomes a degree-$t$ polynomial $q$ satisfying
$q(a)=1$.
The displayed polynomial is feasible, and
$|T_t(x)|\le1$ on $[-1,1]$, so the maximum magnitude of the displayed
polynomial is the right side of
\cref{eq:chebyshev-minimax-value}.

To prove optimality, suppose that another transformed feasible polynomial
$q$ had a
strictly smaller maximum. At the $t+1$ alternating extrema
$x_j=\cos(j\pi/t)$, the polynomial
$T_t(x)-T_t(a)q(x)$ has the sign of $T_t(x_j)=(-1)^j$. The difference
polynomial therefore has at least one root between every consecutive pair of
extrema and also vanishes
at $x=a$ because $q(a)=1$. A nonzero polynomial of degree at most $t$ cannot
have these $t+1$ distinct roots.
\end{proof}

\begin{proposition}[Chebyshev iteration implements the minimax polynomial]
\label{prop:chebyshev-implementation}
Put
\[
 a=\frac{\beta+\mu}{\beta-\mu},
 \qquad C_t=T_t(a).
\]
Starting from $\theta_0$, set
\begin{equation}
 \theta_1
 =\theta_0-\frac{2}{\beta+\mu}\nabla L(\theta_0).
 \label{eq:chebyshev-initial-step}
\end{equation}
For every $t\ge1$, compute
\begin{equation}
 \theta_{t+1}
 =\theta_t
  -\frac{4C_t}{(\beta-\mu)C_{t+1}}\nabla L(\theta_t)
  +\frac{C_{t-1}}{C_{t+1}}(\theta_t-\theta_{t-1}).
 \label{eq:chebyshev-iteration}
\end{equation}
Then
\begin{equation}
 e_t=p_t^\star(H)e_0
 \qquad\text{for every }t\ge0.
 \label{eq:chebyshev-implementation}
\end{equation}
The scalar coefficients can be computed from
\begin{equation}
 C_0=1,\qquad C_1=a,\qquad C_{t+1}=2aC_t-C_{t-1}.
 \label{eq:chebyshev-scalar-recurrence}
\end{equation}
Thus each iteration uses one gradient evaluation, the current iterate, the
preceding iterate, and scalar coefficients determined by $\mu$ and $\beta$.
\end{proposition}

\begin{proof}
The recurrence $T_{t+1}(x)=2xT_t(x)-T_{t-1}(x)$ and
\cref{eq:chebyshev-polynomial} give
\begin{align*}
 p_{t+1}^\star(\lambda)
 &=\frac{2C_t}{C_{t+1}}
   \left(a-\frac{2\lambda}{\beta-\mu}\right)p_t^\star(\lambda)
   -\frac{C_{t-1}}{C_{t+1}}p_{t-1}^\star(\lambda).
\end{align*}
The scalar recurrence in \cref{eq:chebyshev-scalar-recurrence} implies
\[
 \frac{2aC_t}{C_{t+1}}
 =1+\frac{C_{t-1}}{C_{t+1}}.
\]
Substituting $H$ for $\lambda$, applying the resulting polynomial identity to
$e_0$, and using $\nabla L(\theta_t)=He_t$ gives
\cref{eq:chebyshev-iteration}. Finally,
$p_1^\star(\lambda)=1-2\lambda/(\beta+\mu)$ gives
\cref{eq:chebyshev-initial-step}. Induction proves
\cref{eq:chebyshev-implementation}.
\end{proof}

\begin{corollary}[Euclidean-error bound for the minimax polynomial]
\label{cor:chebyshev-rate}
The minimax value in \cref{eq:chebyshev-minimax-value} satisfies
\begin{equation}
 \max_{\lambda\in[\mu,\beta]}|p_t^\star(\lambda)|
 \le2\left(\frac{\sqrt\kappa-1}{\sqrt\kappa+1}\right)^t.
 \label{eq:chebyshev-rate}
\end{equation}
\end{corollary}

\begin{proof}
Put
\[
 r=\frac{\sqrt\kappa-1}{\sqrt\kappa+1}.
\]
Direct calculation gives
\[
 \frac{\beta+\mu}{\beta-\mu}
 =\frac{\kappa+1}{\kappa-1}
 =\frac{r+r^{-1}}2.
\]
The identity
$T_t((u+u^{-1})/2)=(u^t+u^{-t})/2$ therefore gives
\[
 T_t\left(\frac{\beta+\mu}{\beta-\mu}\right)
 =\frac{r^t+r^{-t}}2.
\]
Taking the reciprocal yields
\[
 \frac{2r^t}{1+r^{2t}}\le2r^t,
\]
which proves \cref{eq:chebyshev-rate}.
\end{proof}

\begin{corollary}[Relation between Chebyshev iteration and heavy ball]
\label{cor:chebyshev-heavy-ball}
Let
\[
 r=\frac{\sqrt\kappa-1}{\sqrt\kappa+1}.
\]
The Chebyshev update \cref{eq:chebyshev-iteration} has the heavy-ball form
\begin{equation}
 \theta_{t+1}
 =\theta_t-\eta_t^{\rm Ch}\nabla L(\theta_t)
  +\gamma_t^{\rm Ch}(\theta_t-\theta_{t-1}),
 \label{eq:chebyshev-heavy-ball-form}
\end{equation}
where
\begin{equation}
 \eta_t^{\rm Ch}=\frac{4C_t}{(\beta-\mu)C_{t+1}},
 \qquad
 \gamma_t^{\rm Ch}=\frac{C_{t-1}}{C_{t+1}}.
 \label{eq:chebyshev-heavy-ball-coefficients}
\end{equation}
The coefficients satisfy
\begin{equation}
 \lim_{t\to\infty}\eta_t^{\rm Ch}
 =\frac{4}{(\sqrt\beta+\sqrt\mu)^2},
 \qquad
 \lim_{t\to\infty}\gamma_t^{\rm Ch}=r^2.
 \label{eq:chebyshev-heavy-ball-limits}
\end{equation}
Thus the limits are the quadratic-optimal Polyak heavy-ball coefficients.
Chebyshev iteration uses the time-varying coefficients in
\cref{eq:chebyshev-heavy-ball-coefficients} to attain the finite-degree
minimax polynomial; fixed-coefficient heavy ball does not attain that
polynomial in general.
\end{corollary}

\begin{proof}
The calculation in the proof of \cref{cor:chebyshev-rate} gives
\[
 C_t=\frac{r^t+r^{-t}}2.
\]
Taking the two ratios in
\cref{eq:chebyshev-heavy-ball-coefficients} gives
\[
 \frac{C_t}{C_{t+1}}\longrightarrow r,
 \qquad
 \frac{C_{t-1}}{C_{t+1}}\longrightarrow r^2.
\]
The identity
\[
 \frac{4r}{\beta-\mu}
 =\frac{4}{(\sqrt\beta+\sqrt\mu)^2}
\]
proves \cref{eq:chebyshev-heavy-ball-limits}. For comparison, fixed-coefficient
heavy ball initialized by $\theta_{-1}=\theta_0$ has characteristic
polynomial $(z+r)^2$ at $\lambda=\beta$ and initial values
$p_0^{\rm HB}(\beta)=1$ and $p_1^{\rm HB}(\beta)=-r(2+r)$. Hence
\begin{equation}
 p_t^{\rm HB}(\beta)=(-r)^t\bigl(1+t(1+r)\bigr).
 \label{eq:heavy-ball-endpoint-polynomial}
\end{equation}
The linear factor in \cref{eq:heavy-ball-endpoint-polynomial} shows directly
that the fixed-coefficient method is not the finite-degree minimax method.
\end{proof}

The square root of $\kappa$ comes from uniform polynomial approximation over
the spectral interval. Proposition~\ref{prop:chebyshev-implementation} shows
that the minimax polynomial is also an algorithm: its three-term recurrence
becomes the time-varying heavy-ball update
\cref{eq:chebyshev-heavy-ball-form}. The implementation does not compute an
eigendecomposition, although it does require the spectral bounds $\mu$ and
$\beta$.

Nesterov acceleration generates a third polynomial sequence. On the
quadratic objective, \cref{eq:accelerated-eigendirection} implies
\begin{equation}
 p_{t+1}^{\rm NAG}(\lambda)
 =\left(1-\frac{\lambda}{\beta}\right)
  \left((1+\gamma)p_t^{\rm NAG}(\lambda)
        -\gamma p_{t-1}^{\rm NAG}(\lambda)\right).
 \label{eq:nag-residual-polynomial}
\end{equation}
with
\begin{equation}
 p_0^{\rm NAG}(\lambda)=1,
 \qquad
 p_1^{\rm NAG}(\lambda)=1-\frac{\lambda}{\beta}.
 \label{eq:nag-residual-initial-values}
\end{equation}

\begin{proposition}[Worst-case Nesterov residual on quadratics]
\label{prop:nag-residual-bound}
Set $\tau=1/\sqrt\kappa$ and initialize $\theta_{-1}=\theta_0$. The residual
polynomials defined by \cref{eq:nag-residual-polynomial} satisfy, for every
$t\ge0$,
\begin{equation}
 \max_{\lambda\in[\mu,\beta]}
 |p_t^{\rm NAG}(\lambda)|
 =\left(1+\frac{t}{\sqrt\kappa}\right)
  \left(1-\frac1{\sqrt\kappa}\right)^t.
 \label{eq:nag-residual-bound}
\end{equation}
\end{proposition}

\begin{proof}
The claim is immediate for $t=0$. For $t\ge1$, put
\[
 x=\sqrt{\frac{1-\lambda/\beta}{1-\tau^2}}\in[0,1].
\]
Let $U_j$ denote the Chebyshev polynomial of the second kind, characterized by
$U_j(\cos\phi)=\sin((j+1)\phi)/\sin\phi$ with the endpoint values defined by
continuity. Substitution in \cref{eq:nag-residual-polynomial}, together with
$\gamma=(1-\tau)/(1+\tau)$, gives
\begin{equation}
 p_t^{\rm NAG}(\lambda)
 =\bigl((1-\tau)x\bigr)^t
  \left(T_t(x)+\tau xU_{t-1}(x)\right).
 \label{eq:nag-chebyshev-representation}
\end{equation}
Indeed, after division by $((1-\tau)x)^t$, the recurrence becomes the
Chebyshev recurrence, and
$p_1^{\rm NAG}(\lambda)=(1-\tau)(1+\tau)x^2$.
The representation at $x=0$ follows by continuity.
For $x\in[-1,1]$,
\[
 |T_t(x)|\le1,
 \qquad
 |U_{t-1}(x)|\le t.
\]
Therefore
\[
 |p_t^{\rm NAG}(\lambda)|
 \le (1-\tau)^t x^t(1+\tau tx)
 \le (1-\tau)^t(1+\tau t).
\]
At $\lambda=\mu$, one has $x=1$, $T_t(1)=1$, and $U_{t-1}(1)=t$.
The upper bound is therefore attained.
\end{proof}

Equations~\eqref{eq:chebyshev-rate} and \eqref{eq:nag-residual-bound} both
give $\sqrt\kappa$ iteration dependence, but they are different finite-time
statements. Chebyshev semi-iteration generates the minimax polynomial
\cref{eq:chebyshev-polynomial}; Nesterov acceleration generates
\cref{eq:nag-chebyshev-representation}. The Euclidean-error comparison in
this subsection is also distinct from the objective-gap potential guarantee
in \cref{thm:accelerated-strongly-convex}.

\section{Take-home messages}
\label{sec:take-home}

\begin{itemize}[leftmargin=*]
\item Backtracking line search avoids advance knowledge of the smoothness
constant by testing the finite-step decrease inequality needed in the proof.
Under the PL inequality and with exact objective values, the number of
gradient evaluations needed to reach a prescribed objective gap increases by
at most the search multiplier $\rho$; the additional cost is objective
evaluations at rejected trial points.
\item For smooth, strongly convex objectives, acceleration reduces the number
of gradient evaluations needed for a fixed relative objective-gap reduction
from order $\kappa$ to order $\sqrt\kappa$. A matching first-order lower bound
makes the square-root dependence optimal. The stated one-step potential proof
uses the lower model supplied by strong convexity. The accelerated update can
be derived by combining the smooth upper step with that lower model and
choosing the coupling coefficient so that the squared-gradient terms cancel.
The fixed momentum coefficient uses $\kappa$. For smooth convex objectives
satisfying the PL inequality, restarting the smooth-convex accelerated method
also gives square-root condition-number dependence; a geometric search over
epoch lengths handles an unknown $\kappa$ when $L^\star$ is available.
\item Preconditioning changes coordinates. The relevant quadratic curvature
is $P^{1/2}HP^{1/2}$, and a structured approximation trades geometric
fidelity for near-linear computation and memory.
\item Scalar online adaptation matches the best constant-stepsize regret
bound in hindsight within a factor $3/2$. Diagonal AdaGrad is no worse at
the level of the scalar and diagonal upper bounds and can improve the scalar
bound when gradient activity is concentrated in a small number of coordinates.
\item RMSProp replaces cumulative memory by exponential squared-gradient
memory. Adam adds an exponential average of the gradients and bias correction.
Both updates remain dependent on a base learning rate and memory parameters.
Exponential forgetting prevents a general AdaGrad-style guarantee.
Counterexamples exist in online and
stochastic convex optimization, and a deterministic batch Adam variant can
approach a nonoptimal limit cycle. Restricted positive theory identifies
nonuniform smoothness and favorable coordinate or block geometry as possible
benefits. On the scalar separable logistic example, exact gradient
normalization improves loss decay from $\Theta(1/t)$ to
$\Theta(e^{-\alpha t})$. The conclusion applies to the scalar example;
comparisons on other objectives require separate analysis.
\item AdamW applies uniform weight decay outside the adaptive gradient
scaling. An $\ell_2$ penalty inside the objective receives coordinatewise
scaling from the adaptive matrix. The two formulas define different updates.
\item Universal deterministic methods can attain the rate associated with an
unknown H\"older-smoothness regime. In stochastic convex optimization, methods
can match tuned SGD up to logarithmic factors on bounded domains. The analogous
guarantee over an unbounded domain is impossible under the general
bounded-noise oracle model. Every tuning-free claim must specify the retained
hints, oracle, domain, reference method, and performance measure.
\end{itemize}

\printendnotes

\section*{Bibliographical notes}

The sufficient-decrease line search goes back to classical work including
\citet{armijo1966}. The presentation here uses the quadratic upper-bound test
that appears in proximal-gradient and accelerated proximal-gradient methods.
\citet{beck2009} analyze both ordinary and accelerated backtracking and prove
the $O(1/t)$ and $O(1/t^2)$ convex objective-gap bounds used here.

Polyak's heavy-ball method was introduced by \citet{polyak1964}. Its
quadratic-optimal parameters give accelerated asymptotic spectral convergence
on strongly convex quadratics. Nesterov's accelerated method was introduced
by \citet{nesterov1983}; \citet{nesterov2004} gives a systematic treatment of
its guarantees over the full smooth convex and smooth strongly convex
classes, together with the standard first-order lower bounds. The construction
in \cref{sec:acceleration} is an elementary estimate-sequence argument: it
combines a smooth upper step with accumulated strongly convex lower models and
determines the coupling coefficient from the desired potential contraction.
The smooth-convex construction in
\cref{sec:convex-acceleration-lyapunov} expresses the estimate-sequence
mechanism as a time-varying Lyapunov inequality.

Restarting converts the smooth-convex accelerated rate into a linear rate
under quadratic growth. The implication from PL to quadratic growth, with the
same constant, is proved by \citet{karimi2016}.
\citet{odonoghue2015} propose objective- and gradient-based adaptive restart
rules motivated by oscillation of the
accelerated iterates. \citet{fercoq2019} give a provable adaptive restart
scheme under a local quadratic-growth condition; their test uses the norm of
the gradient mapping and adjusts the restart frequency from observed
progress. \citet{roulet2020} analyze scheduled and adaptive restarts under
more general H\"olderian growth conditions. For convex objectives, such growth
conditions are related to the nonlinear gradient-dominance inequalities in
Lecture~11 and lead to exponent-dependent restart schedules and rates. The
main text uses the quadratic-growth case because it gives the clean convex-PL
result in \cref{eq:accelerated-restart-main-complexity}.

The distinction between quadratic and general objectives is essential for
heavy ball. \citet{lessard2016} give a smooth, strongly convex nonquadratic
example on which the quadratic-optimal heavy-ball parameters produce a
nonconvergent cycle. In dimension at least two, \citet{goujaud2023} prove a
stronger fixed-parameter non-acceleration result: every fixed heavy-ball
parameter choice either has a nonaccelerated worst-case rate on strongly
convex quadratics or fails to converge on some smooth, strongly convex
objective.

The polynomial construction in \cref{sec:chebyshev} is the classical
Chebyshev semi-iterative method. \citet{barrett1994} present the three-term
implementation for linear systems and discuss its dependence on spectral
bounds. Chebyshev semi-iteration is a nonstationary heavy-ball iteration: its
coefficients converge to the quadratic-optimal Polyak coefficients, but the
time-varying coefficients are what generate the finite-degree minimax
polynomial. Nesterov acceleration generates a different polynomial recurrence.
All three methods exhibit square-root condition-number dependence in their
respective quadratic guarantees. The lower bound separating smooth PL
objectives from strongly convex objectives is proved by \citet{yue2023}: over
the general smooth PL class, first-order methods require order
$(\beta/\mu)\log(1/\epsilon)$ oracle calls in the worst case.

The Kronecker-factored calculation in
\cref{eq:exact-fisher-block,eq:kfac-block} is the structural approximation
used by K-FAC. \citet{zhang2019ngd} prove global convergence for natural
gradient and K-FAC under Jacobian rank and stability conditions, including an
overparameterized two-layer ReLU setting. \citet{karakida2020} analyze deep
networks in an infinite-width limit and show that several approximate Fisher
matrices have the same function-space dynamics as exact natural gradient
under the assumptions of their analysis. The two sets of results justify
specific structured approximations in restricted regimes; the results do not
give a general finite-width conditioning theorem for multilayer networks.

The adaptive subgradient framework and data-dependent regret guarantees for
AdaGrad are due to \citet{duchi2011}. The diagonal theorem in
\cref{thm:adagrad-regret} is a direct coordinatewise specialization with a
simple box constraint. \Cref{ex:l12-adagrad-regret} asks you to prove the
specialization and identify exactly which observed geometry enters.

RMSProp was introduced in Geoffrey Hinton's 2012 neural-network course
\citep{hinton2012rmsprop}. Its update has several implementation variants,
including placement of
$\varepsilon$ and optional correction of the initial second-moment bias.
\citet{zou2019} give sufficient conditions for convergence of RMSProp and
Adam in nonconvex stochastic optimization. The conditions involve the base
stepsizes and the accumulated second-moment quantities. The result supplies a
conditional convergence guarantee. A tuning-free comparison with SGD requires
a different theorem.

Adam was introduced by \citet{kingma2015}. The original convergence analysis
was later shown to contain a problem. \citet{reddi2018} prove separate online
and stochastic convex nonconvergence results and propose AMSGrad. The
complete-forgetting example in \cref{ex:adam-forgetting} is an elementary
limiting case that isolates the online mechanism. \citet{bock2022} construct
trajectories of a deterministic bias-corrected batch Adam variant that
approach a nonoptimal period-two limit cycle on a scalar strongly convex
quadratic. Their denominator uses $\sqrt{v_t+\varepsilon}$, and the result
rules out global convergence without asserting failure from every
initialization. \citet{defossez2022}
give a convergence analysis for Adam and AdaGrad on smooth, possibly
nonconvex objectives with bounded gradients. Results aimed at explaining an
advantage beyond basic convergence include the nonuniform-smoothness
comparison of \citet{wang2024}, the $\ell_\infty$ and blockwise analysis of
\citet{xie2025}, and the block-heterogeneity analysis and experiments of
\citet{zhang2024transformers}. The assumptions and conclusions of the three
analyses differ, so none provides a universal ordering of Adam and SGD.

The original motivation for Adam combines AdaGrad's behavior on sparse
gradients with RMSProp's responsiveness to nonstationary gradient scales; see
\citet{kingma2015}. The original paper also reports experiments on
$\ell_2$-regularized multiclass logistic regression and sparse bag-of-words
classification. In the logistic experiment, Adam and momentum SGD have
similar convergence and both improve on AdaGrad under the reported tuning.
The scalar calculation in
\cref{ex:logistic-normalization} isolates the flattening mechanism without
attributing every feature of Adam to that example. For separable logistic
regression, \citet{soudry2018} prove $1/t$ loss decay and convergence toward
the Euclidean max-margin direction for gradient descent with an admissible
constant stepsize.
\citet{zhang2024adambias} prove that full-batch Adam with a class of
diminishing learning rates instead approaches an $\ell_\infty$ max-margin
direction. The latter result concerns implicit bias. Loss-decay rates require
a separate analysis. \citet{axiotis2023} obtain linear logistic-loss
convergence with a variable gradient-descent stepsize, showing that the faster
rate can also arise from gradient descent. Recent nonuniform-smoothness results give sharper
comparisons for normalized methods and restricted Adam variants; see
\citet{vaswani2026}.

Decoupled weight decay and AdamW are due to \citet{loshchilov2019}.
\citet{zhuang2022} interpret AdamW through an approximate proximal update and
study the scale-freeness of the update. Convergence-rate results for AdamW are
also available under explicit smoothness and noise assumptions; see, for
example, \citet{li2025adamw}. The main text retains only the algebraic
distinction between an adaptively scaled $\ell_2$ penalty and decoupled decay
because the distinction is needed to understand the update.

Momentum methods have also been used directly in neural-network training.
\citet{sutskever2013} study classical and Nesterov-style momentum for deep and
recurrent networks, while the ResNet training recipe of \citet{he2016resnet}
uses momentum SGD. \citet{loshchilov2017sgdr} introduce cosine annealing with
warm restarts. A modern large-language-model example is the Llama~2 recipe of
\citet{touvron2023llama2}, which combines AdamW, learning-rate warmup, and
cosine decay. These references document successful uses of the methods; they
do not extend the smooth strongly convex acceleration theorem to neural
networks or establish a universal ordering of optimizers.

The universal accelerated method of \citet{nesterov2015} is given the target
accuracy while the H\"older exponent and its constant remain unknown. Its
complexity gives both endpoint rates in
\cref{eq:universal-convex-complexity}. \Citet{khaled2024} define a tuning-free
counterpart of SGD by comparison with optimally tuned SGD under loose
parameter hints. They identify existing positive results on bounded domains
and prove the unbounded-domain impossibility summarized in
\cref{prop:tuning-free-scope}. They also give positive unbounded-domain
results under stronger noise assumptions and treat a separate nonconvex
stationarity question. These results ask for simultaneous adaptation to
smoothness, noise, and distance scales, which is a stronger requirement than
backtracking with exact objective evaluations.

\section*{Glossary}

\begin{description}[style=nextline,leftmargin=1.5em,labelindent=0pt]
\item[Acceptable curvature estimate]
A number $M>0$ for which the proposed point $\theta^+(M)$ satisfies the
finite-step acceptance test used by a line search.
\item[Acceleration]
An improvement in worst-case iteration complexity obtained from state carried
across iterations. For $\beta$-smooth, $\mu$-strongly convex objectives, the
gradient complexities are $O(\kappa\log(1/\epsilon))$ for gradient descent and
$O(\sqrt\kappa\log(1/\epsilon))$ for accelerated gradient descent. For
$\beta$-smooth convex objectives with no PL or strong-convexity assumption,
the corresponding objective-gap bounds are $O(\beta R^2/t)$ and
$O(\beta R^2/t^2)$, where $R=\|\theta_0-\theta^\star\|_2$.
\item[Adam]
An adaptive stochastic first-order method that keeps exponential averages
$m_t$ and $v_t$ of the gradients and squared gradients. The update
\cref{eq:adam-update} divides the bias-corrected first-moment estimate
$\widehat m_t$ coordinatewise by $\sqrt{\widehat v_t}+\varepsilon$.
\item[AdamW and decoupled weight decay]
An Adam update in which the parameter shrinkage is applied separately from
the adaptively scaled gradient:
$\theta_{t+1}=(1-\alpha\lambda)\theta_t-\alpha D_t^{-1}g_t$.
\item[Backtracking line search]
A procedure that increases a trial curvature estimate, equivalently decreases
the trial stepsize, until an evaluated finite-step inequality holds.
Accelerated backtracking also updates the coupling coefficients that depend on
the accepted curvature estimate.
\item[Chebyshev semi-iteration]
A quadratic-optimization method whose time-varying two-step recurrence
implements the degree-$t$ polynomial minimizing
$\max_{\lambda\in[\mu,\beta]}|p_t(\lambda)|$ subject to $p_t(0)=1$.
\item[Cosine learning-rate decay]
A fixed-budget schedule that changes the global learning rate according to
\cref{eq:cosine-learning-rate}, from $\alpha_{\max}$ at the start of the
decay to $\alpha_{\min}$ at its end.
\item[Diagonal AdaGrad]
An online first-order method whose coordinate-$i$ denominator is the square
root of the cumulative squared gradients observed in that coordinate.
\item[Exponential memory]
A moving average in which an observation from $s$ iterations ago receives a
weight proportional to a power such as $\beta_2^s$.
\item[First-order oracle]
An interface that returns an objective value and gradient at a queried
parameter. In the lower-bound comparison in this lecture, each query may
depend on all preceding queries and oracle replies.
\item[Heavy-ball method]
The two-step update
$\theta_{t+1}=\theta_t-\eta\nabla L(\theta_t)
+\gamma(\theta_t-\theta_{t-1})$, which evaluates the gradient at the current
iterate. Quadratic-optimal fixed coefficients need not converge on every
smooth, strongly convex objective.
\item[Hindsight-tuned stepsize]
A constant stepsize chosen after the complete gradient sequence is known to
minimize a proved upper bound such as \cref{eq:fixed-eta-oco}. The comparison
concerns the upper bound, not necessarily the realized losses.
\item[Nesterov accelerated gradient]
The two-step method
$y_t=\theta_t+\gamma(\theta_t-\theta_{t-1})$ and
$\theta_{t+1}=y_t-\nabla L(y_t)/\beta$. For the strongly convex theorem in
this lecture, $\gamma=(\sqrt\kappa-1)/(\sqrt\kappa+1)$.
\item[Online convex optimization and regret]
At round $t$, the learner chooses $x_t$ before observing the convex loss
$\ell_t$. Its regret relative to a fixed $u$ is
$\operatorname{Reg}_T(u)=\sum_{t=1}^T(\ell_t(x_t)-\ell_t(u))$.
\item[Preconditioner]
A positive-definite matrix $P$ multiplying the gradient. On a quadratic with
Hessian $H$, the conditioning effect of $P$ is determined by
$P^{1/2}HP^{1/2}$.
\item[Quadratic growth]
For a nonempty minimizer set $\Theta^\star$, the inequality
$L(\theta)-L^\star\ge
\mu\operatorname{dist}(\theta,\Theta^\star)^2/2$, where
$\operatorname{dist}(\theta,\Theta^\star)
=\inf_{z\in\Theta^\star}\|\theta-z\|_2$. The PL inequality implies quadratic
growth with the same constant. In the convex setting, this bound converts the
distance-dependent accelerated rate into a gap-dependent restart guarantee.
\item[Restart]
Reinitialization of an accelerated method's memory after a fixed or
data-dependent epoch while retaining the current parameter. Restarting a
smooth-convex accelerated method every $\Theta(\sqrt\kappa)$ iterations gives
linear convergence under convexity and the PL inequality; adaptive restart
rules search for this epoch length when $\kappa$ is unknown.
\item[Scalar online adaptation]
A stepsize rule based on the cumulative Euclidean gradient magnitude. The
rule in \cref{thm:scalar-adaptation} matches the best constant-stepsize bound
in hindsight within a factor $3/2$ without knowing the horizon or future
gradient magnitudes.
\item[RMSProp]
An adaptive stochastic first-order method that divides the current gradient
coordinatewise by the square root of an exponential average of recent squared
gradients, as in \cref{eq:rmsprop-update}.
\item[Tuning-free guarantee]
A theorem comparing an algorithm that lacks specified problem parameters
with a benchmark that is optimally tuned using those parameters. The omitted
parameters, allowed hints, oracle, and domain are part of the definition.
\item[Universal convex method]
A method that attains the rate associated with the unknown H\"older parameters
$\nu$ and $H_\nu$ in \cref{eq:holder-first-order-oracle}. The guarantee
\cref{eq:universal-convex-complexity} includes the nonsmooth and
smooth-convex endpoint rates.
\item[Warmup phase]
An initial part of a learning-rate schedule during which the global learning
rate increases from a smaller value to its prescribed peak value.
\end{description}

\section{Exercises}
\label{sec:exercises}

\exerciseinstructions
\setcounter{courseexercise}{0}
\exercisegroup{Unknown smoothness and acceleration}
\begin{exercisechunk}
\item \textbf{Backtracking avoids conservative curvature scales.}
\label[exercise]{ex:l12-backtracking-quadratic}
\exerciseprofile{2}{2}{3}{\exproof\quad\excalculation\quad\exinterpretation}{%
Quantify the work saved by testing a locally sufficient curvature scale.}
Let
\[
 L(\theta)=\frac{\beta}{2}\theta^2,
 \qquad \beta>0,
\]
and consider the trial point
$\theta^+(M)=\theta-L'(\theta)/M$ for $M>0$.
The sufficient-decrease test is
\begin{equation}
 L(\theta^+(M))\le L(\theta)-\frac{1}{2M}|L'(\theta)|^2.
 \label{eq:ex-l12-backtracking-test}
\end{equation}
\begin{enumerate}[(a),beginpenalty=10000]
\item For $\theta\ne0$, prove that \cref{eq:ex-l12-backtracking-test} holds if
and only if $M\ge\beta$. State what happens when $\theta=0$.

\item Let $0<a<b$ and define
\[
 \widetilde L(\theta)=
 \begin{cases}
  a\theta^2/2,&\theta\le0,\\
  b\theta^2/2,&\theta>0.
 \end{cases}
\]
Prove that $\widetilde L$ is differentiable and $b$-smooth. Starting from
$\theta_0<0$,
show that the globally certified step $1/b$ gives
\[
 \theta_t=\left(1-\frac ab\right)^t \theta_0,
\]
whereas the trial curvature $M=a$ passes
\cref{eq:ex-l12-backtracking-test} with $\widetilde L$ in place of $L$ and reaches the
minimizer in one step. Deduce that the globally certified method
requires
\[
 t\ge
 \frac{\log(1/\epsilon)}{-\log(1-a/b)}
\]
iterations to obtain $|\theta_t|\le\epsilon|\theta_0|$, for
$0<\epsilon<1$.

\item Return to $L(\theta)=\beta\theta^2/2$. Starting from an arbitrary
$M_0>0$, repeatedly replace $M$ by $2M$ until the test holds. Prove that the accepted
value satisfies
\[
 M_{\rm acc}\le\max\{M_0,2\beta\}
\]
and that the number of rejected trials is at most
\[
 \max\left\{0,\left\lceil\log_2\frac{\beta}{M_0}\right\rceil\right\}.
\]
\end{enumerate}
\begin{hints}
For (a), expand $(1-\beta/M)^2$. For (b), check pairs of points on the same
side of zero and pairs that straddle zero separately.
\end{hints}
\end{exercisechunk}

\begin{exercisechunk}
\item \textbf{Eliminating the auxiliary state and proving the accelerated rate.}
\label[exercise]{ex:l12-nesterov-equivalence}
\exerciseprofile{2}{2}{2}{\exproof\quad\excalculation}{%
Eliminate the auxiliary state and use the potential contraction to prove the
accelerated linear rate.}
Let $L:\R^d\to\R$ be $\beta$-smooth and $\mu$-strongly convex, let
$\theta^\star$ be its minimizer, and define
\[
 \kappa=\frac\beta\mu,
 \qquad
 \tau=\frac1{\sqrt\kappa},
 \qquad
 \gamma=\frac{1-\tau}{1+\tau}.
\]
Initialize $v_0=\theta_0$ and $\theta_{-1}=\theta_0$.
\begin{enumerate}[(a),beginpenalty=10000]
\item Prove that the coupled update
\begin{align*}
 y_t&=\frac{\theta_t+\tau v_t}{1+\tau},
 &\theta_{t+1}&=y_t-\frac1\beta\nabla L(y_t),\\
 v_{t+1}&=(1-\tau)v_t
 +\tau\left(y_t-\frac1\mu\nabla L(y_t)\right)
\end{align*}
and the two-step update
\[
 y_t=\theta_t+\gamma(\theta_t-\theta_{t-1}),
 \qquad
 \theta_{t+1}=y_t-\frac1\beta\nabla L(y_t)
\]
generate the same sequences $(\theta_t)$ and $(y_t)$.
\item Suppose the coupled update satisfies
\[
 V(\theta_t,v_t)\le(1-\tau)^tV(\theta_0,v_0),
 \qquad
 V(\theta,v)=L(\theta)-L^\star
 +\frac\mu2\|v-\theta^\star\|_2^2.
\]
Prove
\[
 L(\theta_t)-L^\star
 \le2\left(1-\frac1{\sqrt\kappa}\right)^t
       \bigl(L(\theta_0)-L^\star\bigr),
\]
and, for $0<\epsilon<2(L(\theta_0)-L^\star)$, show that
$L(\theta_t)-L^\star\le\epsilon$ whenever
\[
 t\ge\sqrt\kappa
 \log\frac{2(L(\theta_0)-L^\star)}{\epsilon}.
\]
\end{enumerate}
\begin{hints}
For part (a), use induction with
\[
 v_t=\frac{\theta_t-(1-\tau)\theta_{t-1}}{\tau}.
\]
For part (b), use strong convexity at $\theta_0$ and
$1-\tau\le e^{-\tau}$.
\end{hints}
\end{exercisechunk}

\begin{exercisechunk}
\item \textbf{Restarting acceleration under convex PL.}
\label[exercise]{ex:l12-accelerated-restart}
\exerciseprofile{2}{2}{3}{\exproof\quad\excalculation\quad\exinterpretation}{%
Convert the accelerated convex rate into square-root condition-number
dependence and adapt the restart length by geometric search.}
Let $L:\R^d\to\R$ be $\beta$-smooth, convex, and $\mu$-PL, let
$\kappa=\beta/\mu$, and suppose that its minimizer set $\Theta^\star$ is
nonempty. You may use the quadratic-growth consequence of the PL inequality,
\begin{equation}
 \operatorname{dist}(z,\Theta^\star)^2
 \le\frac2\mu\bigl(L(z)-L^\star\bigr).
 \label{eq:ex-l12-pl-quadratic-growth}
\end{equation}
Here
$\operatorname{dist}(z,\Theta^\star)
=\inf_{\theta^\star\in\Theta^\star}\|z-\theta^\star\|_2$.
Suppose that an
accelerated method started from any $z$ produces, after $s\ge1$ iterations,
an output $\mathsf A_s(z)$ satisfying
\begin{equation}
 L(\mathsf A_s(z))-L^\star
 \le \frac{2\beta\operatorname{dist}(z,\Theta^\star)^2}{(s+1)^2}.
 \label{eq:ex-l12-accelerated-rate}
\end{equation}
\begin{enumerate}[(a),beginpenalty=10000]
\item Prove
\[
 L(\mathsf A_s(z))-L^\star
 \le \frac{4\kappa}{(s+1)^2}\bigl(L(z)-L^\star\bigr).
\]
\item Choose an integer epoch length $s$ that makes the multiplicative
factor at most $1/2$. Starting each epoch from the preceding epoch's output,
prove that $r$ epochs give objective gap at most
$2^{-r}(L(\theta_0)-L^\star)$.
\item Derive a total iteration bound of order
\[
 \sqrt\kappa\log\frac{L(\theta_0)-L^\star}{\epsilon}
\]
for reaching objective gap at most $\epsilon$. State which of the two
inequalities used in part (a) comes from smoothness and convexity and which
comes from the PL inequality.
\item Suppose that $\kappa$ is unknown but $L^\star$ is available. From a
fixed point $z$, try epoch lengths $s=1,2,4,\ldots$, restarting every trial
from $z$, until
\[
 L(\mathsf A_s(z))-L^\star
 \le\frac12\bigl(L(z)-L^\star\bigr).
\]
Prove that the first successful trial has $s<2\sqrt{8\kappa}$ and that all
trials together use fewer than $2s$ gradient evaluations. If the successful
epoch length is then reused, with each later epoch starting from the preceding
epoch's output, derive the total bound
\[
 O\left(\sqrt\kappa\left(
 1+\log\frac{L(\theta_0)-L^\star}{\epsilon}\right)\right).
\]
\end{enumerate}
\begin{hints}
For part (b), require $(s+1)^2\ge8\kappa$.
\end{hints}
\end{exercisechunk}

\exercisegroup{Adaptive geometry}
\begin{exercisechunk}
\item \textbf{A rotated quadratic defeats diagonal preconditioning.}
\label[exercise]{ex:l12-rotated-preconditioner}
\exerciseprofile{2}{2}{2}{\exproof\quad\excalculation\quad\exinterpretation}{%
Distinguish coordinatewise rescaling from correction of rotated curvature.}
Let $\kappa>1$ and
\[
 Q=\frac1{\sqrt2}\begin{pmatrix}1&-1\\1&1\end{pmatrix},
 \qquad
 H=Q\begin{pmatrix}1&0\\0&\kappa\end{pmatrix}Q^\top,
 \qquad
 L(\theta)=\frac12\theta^\top H\theta.
\]
For a positive-definite preconditioner $P$, the transformed Hessian is
$P^{1/2}HP^{1/2}$.
\begin{enumerate}[(a),beginpenalty=10000]
\item Compute $H$ and prove that its condition number is $\kappa$.
\item Let $P_{\rm diag}=\diag(H)^{-1}$. Compute
$P_{\rm diag}^{1/2}HP_{\rm diag}^{1/2}$ and prove that its condition number
remains $\kappa$.
\item Let $P_{\rm full}=H^{-1}$. Compute the transformed Hessian and its
condition number. Explain precisely which geometric information is absent
from $P_{\rm diag}$.
\end{enumerate}
\begin{hints}
The two diagonal entries of $H$ are equal. Orthogonal similarity preserves
eigenvalues.
\end{hints}
\end{exercisechunk}

\begin{exercisechunk}
\item \textbf{The diagonal AdaGrad bound.}\suppmark
\label[exercise]{ex:l12-adagrad-regret}
\exerciseprofile{3}{3}{3}{\exproof\quad\excalculation\quad\exinterpretation}{%
Derive a regret bound whose coordinate scales are determined by the observed
gradients.}
Let $\cX=\prod_{i=1}^d[a_i,b_i]$ with
$D_i=b_i-a_i>0$. On round $t$, let $\ell_t:\cX\to\R$ be convex and let
$g_t\in\partial\ell_t(x_t)$. Define
\[
 G_{t,i}=\left(\sum_{s=1}^t g_{s,i}^2\right)^{1/2},
 \qquad
 \eta_{t,i}=\frac{D_i}{G_{t,i}}
\]
when $G_{t,i}>0$. If $G_{t,i}=0$, leave coordinate $i$ unchanged. Otherwise
update it by
\[
 x_{t+1,i}=\Pi_{[a_i,b_i]}(x_{t,i}-\eta_{t,i}g_{t,i}).
\]
\begin{enumerate}[(a),beginpenalty=10000]
\item For nonnegative $a_1,\ldots,a_T$ and
$A_t=\sum_{s=1}^ta_s$, prove
\[
 \sum_{t=1}^T\frac{a_t}{\sqrt{A_t}}\le2\sqrt{A_T},
\]
where a term with $A_t=0$ is defined to be zero.
\item Fix $u\in\cX$. Prove for each coordinate that
\[
 \sum_{t=1}^T g_{t,i}(x_{t,i}-u_i)
 \le\frac32D_iG_{T,i}.
\]
\item Use convexity to prove the regret bound
\[
 \sum_{t=1}^T\bigl(\ell_t(x_t)-\ell_t(u)\bigr)
 \le\frac32\sum_{i=1}^dD_i
       \left(\sum_{t=1}^Tg_{t,i}^2\right)^{1/2}.
\]
State how the right side differs from a bound based only on one worst-case
gradient norm.
\end{enumerate}
\begin{hints}
For (a), use
$a_t/\sqrt{A_t}\le2(\sqrt{A_t}-\sqrt{A_{t-1}})$.
For (b), use nonexpansiveness of interval projection and telescope with the
nonincreasing stepsizes.
\end{hints}
\end{exercisechunk}

\exercisegroup{Changing gradient scales and weight decay}
\begin{exercisechunk}
\item \textbf{Complete forgetting can give linear regret.}
\label[exercise]{ex:l12-adam-forgetting}
\exerciseprofile{2}{2}{2}{\exproof\quad\excalculation\quad\exinterpretation}{%
Verify a convex failure mechanism caused by normalizing only with the current
gradient magnitude.}
Let $\cX=[-1,1]$, choose $C>2$ and $\alpha\in(0,2]$, and repeat the three
linear losses
\[
 \ell_{3k+1}(x)=Cx,
 \qquad
 \ell_{3k+2}(x)=\ell_{3k+3}(x)=-x,
 \qquad k\ge0.
\]
Consider the projected Adam special case with no momentum, complete
second-moment forgetting, and no stabilizer:
\[
 m_t=g_t,
 \qquad v_t=g_t^2,
 \qquad
 x_{t+1}=\Pi_{[-1,1]}
 \left(x_t-\alpha\frac{m_t}{\sqrt{v_t}}\right),
 \qquad x_1=1,
\]
where $g_t=\ell_t'(x_t)$.
\begin{enumerate}[(a),beginpenalty=10000]
\item Prove that every three-round block begins at $x_{3k+1}=1$ and list the
three iterates used to evaluate its losses.
\item Prove that the minimizer of the sum of the three losses over $[-1,1]$
is $u=-1$. Compute the three-round regret against $u$ and show that it equals
$2C-4+\alpha>0$.
\item Conclude that the regret after $3K$ rounds grows linearly in $K$.
Explain why the squared-gradient statistic in this example is not a reliable
curvature estimate.
\end{enumerate}
\begin{hints}
The normalized update uses only $\operatorname{sign}(g_t)$. Track the
projection after each of the two gradients equal to $-1$.
\end{hints}
\end{exercisechunk}

\begin{exercisechunk}
\item \textbf{Adaptive penalties and decoupled weight decay.}
\label[exercise]{ex:l12-adamw-decoupling}
\exerciseprofile{1}{1}{1}{\excalculation\quad\exinterpretation}{%
Distinguish an adaptively preconditioned $\ell_2$ penalty from uniform weight
decay.}
Let $D_t$ be a positive diagonal matrix, let $g_t=\nabla L(x_t)$, and let
$\alpha,\lambda>0$.
\begin{enumerate}[(a),beginpenalty=10000]
\item Apply the adaptive update $x_{t+1}=x_t-\alpha D_t^{-1}\nabla F(x_t)$
to
\[
 F(x)=L(x)+\frac\lambda2\|x\|_2^2.
\]
Show that
\[
 x_{t+1}=(I-\alpha\lambda D_t^{-1})x_t
          -\alpha D_t^{-1}g_t.
\]
\item AdamW instead uses
\[
 x_{t+1}=(1-\alpha\lambda)x_t-\alpha D_t^{-1}g_t.
\]
State the necessary and sufficient condition on $D_t$ under which the two
shrinkage maps coincide for every $x_t$. Explain what is decoupled in the
AdamW update.
\end{enumerate}
\begin{hints}
Differentiate the quadratic penalty. Compare the two matrices multiplying
$x_t$.
\end{hints}
\end{exercisechunk}

\exercisegroup{Comparing adaptive and tuned guarantees}
\begin{exercisechunk}
\item \textbf{Comparing adaptive methods with tuned methods.}\suppmark
\label[exercise]{ex:l12-tuned-comparisons}
\exerciseprofile{2}{2}{3}{\excalculation\quad\exinterpretation}{%
Compare the guarantees proved in the lecture with the corresponding bounds
obtained by tuning the stepsizes in advance.}
\begin{enumerate}[(a),beginpenalty=10000]
\item Suppose $M_{\mathrm{init}}\le\beta$. Compare
\cref{thm:backtracking-pl,eq:backtracking-rejections} with gradient descent
using the tuned stepsize $1/\beta$. Show that the guaranteed contraction factor
changes from $1-\mu/\beta$ to $1-\mu/(\rho\beta)$ and that the search uses at
most
\[
 \left\lceil\log_\rho\frac{\beta}{M_{\mathrm{init}}}\right\rceil
\]
additional objective evaluations.

\item Define
\[
 S_T=\left(\sum_{t=1}^T\|g_t\|_2^2\right)^{1/2},
 \qquad
 S_i=\left(\sum_{t=1}^Tg_{t,i}^2\right)^{1/2}.
\]
Prove, for $a,b>0$, that
\[
 \inf_{\eta>0}\left\{\frac{a^2}{2\eta}+\frac{\eta b^2}{2}\right\}
 =ab,
 \qquad \eta^\star=\frac ab.
\]
Apply the identity to \cref{eq:fixed-eta-oco} and coordinatewise to prove
\[
 \inf_{\eta>0}
 \left\{\frac{D^2}{2\eta}+\frac{\eta S_T^2}{2}\right\}=DS_T,
 \qquad
 \inf_{\eta_1,\ldots,\eta_d>0}
 \sum_{i=1}^d\left(\frac{D_i^2}{2\eta_i}
                    +\frac{\eta_iS_i^2}{2}\right)
 =\sum_{i=1}^dD_iS_i.
\]
Use \cref{thm:scalar-adaptation,thm:adagrad-regret} to show that each adaptive
upper bound is at most $3/2$ times its corresponding hindsight-tuned upper
bound. Interpret a term with $S_i=0$ as an infimum over $\eta_i>0$.
\end{enumerate}
\end{exercisechunk}


\clearpage
\begin{thebibliography}{99}

\bibitem[Armijo(1966)]{armijo1966}
Larry Armijo.
\newblock Minimization of functions having Lipschitz continuous first partial
derivatives.
\newblock \emph{Pacific Journal of Mathematics}, 16(1):1--3, 1966.
\newblock \url{https://projecteuclid.org/euclid.pjm/1102995080}.

\bibitem[Axiotis and Sviridenko(2023)]{axiotis2023}
Kyriakos Axiotis and Maxim Sviridenko.
\newblock Gradient descent converges linearly for logistic regression on
separable data.
\newblock In \emph{Proceedings of the 40th International Conference on
Machine Learning}, pages 1302--1319, 2023.
\newblock \url{https://proceedings.mlr.press/v202/axiotis23a.html}.

\bibitem[Barrett et al.(1994)]{barrett1994}
Richard Barrett, Michael Berry, Tony F. Chan, James Demmel, June Donato,
Jack Dongarra, Victor Eijkhout, Roldan Pozo, Charles Romine, and Henk van der
Vorst.
\newblock \emph{Templates for the Solution of Linear Systems: Building Blocks
for Iterative Methods}.
\newblock SIAM, second edition, 1994.
\newblock \url{https://www.netlib.org/templates/templates.pdf}.

\bibitem[Beck and Teboulle(2009)]{beck2009}
Amir Beck and Marc Teboulle.
\newblock A fast iterative shrinkage-thresholding algorithm for linear inverse
problems.
\newblock \emph{SIAM Journal on Imaging Sciences}, 2(1):183--202, 2009.
\newblock \url{https://doi.org/10.1137/080716542}.

\bibitem[Bock and Wei\ss(2022)]{bock2022}
Sebastian Bock and Martin Georg Wei\ss.
\newblock Non-convergence and limit cycles in the Adam optimizer.
\newblock 2022.
\newblock \url{https://arxiv.org/abs/2210.02070}.

\bibitem[D\'efossez et al.(2022)]{defossez2022}
Alexandre D\'efossez, L\'eon Bottou, Francis Bach, and Nicolas Usunier.
\newblock A simple convergence proof of Adam and AdaGrad.
\newblock \emph{Transactions on Machine Learning Research}, 2022.
\newblock \url{https://arxiv.org/abs/2003.02395}.

\bibitem[Duchi et al.(2011)]{duchi2011}
John Duchi, Elad Hazan, and Yoram Singer.
\newblock Adaptive subgradient methods for online learning and stochastic
optimization.
\newblock \emph{Journal of Machine Learning Research}, 12:2121--2159, 2011.
\newblock \url{https://www.jmlr.org/papers/v12/duchi11a.html}.

\bibitem[Fercoq and Qu(2019)]{fercoq2019}
Olivier Fercoq and Zheng Qu.
\newblock Adaptive restart of accelerated gradient methods under local
quadratic growth condition.
\newblock \emph{IMA Journal of Numerical Analysis}, 39(4):2069--2095, 2019.
\newblock \url{https://doi.org/10.1093/imanum/drz007}.

\bibitem[Goujaud et al.(2023)]{goujaud2023}
Baptiste Goujaud, Adrien B. Taylor, and Aymeric Dieuleveut.
\newblock Provable non-accelerations of the heavy-ball method.
\newblock 2023.
\newblock \url{https://arxiv.org/abs/2307.11291}.

\bibitem[He et al.(2016)]{he2016resnet}
Kaiming He, Xiangyu Zhang, Shaoqing Ren, and Jian Sun.
\newblock Deep residual learning for image recognition.
\newblock In \emph{Proceedings of the IEEE Conference on Computer Vision and
Pattern Recognition}, pages 770--778, 2016.
\newblock \url{https://openaccess.thecvf.com/content_cvpr_2016/html/He_Deep_Residual_Learning_CVPR_2016_paper.html}.

\bibitem[Hinton(2012)]{hinton2012rmsprop}
Geoffrey Hinton.
\newblock RMSProp: Divide the gradient by a running average of its recent
magnitude.
\newblock Lecture 6e, \emph{Neural Networks for Machine Learning}, 2012.
\newblock \url{https://www.cs.utoronto.ca/~hinton/coursera_lectures.html}.

\bibitem[Karakida and Osawa(2020)]{karakida2020}
Ryo Karakida and Kazuki Osawa.
\newblock Understanding approximate Fisher information for fast convergence
of natural gradient descent in wide neural networks.
\newblock In \emph{Advances in Neural Information Processing Systems 33},
2020.
\newblock \url{https://papers.neurips.cc/paper_files/paper/2020/hash/7b41bfa5085806dfa24b8c9de0ce567f-Abstract.html}.

\bibitem[Karimi et al.(2016)]{karimi2016}
Hamed Karimi, Julie Nutini, and Mark Schmidt.
\newblock Linear convergence of gradient and proximal-gradient methods under
the Polyak--\L{}ojasiewicz condition.
\newblock In \emph{Machine Learning and Knowledge Discovery in Databases},
pages 795--811, 2016.
\newblock \url{https://arxiv.org/abs/1608.04636}.

\bibitem[Khaled and Jin(2024)]{khaled2024}
Ahmed Khaled and Chi Jin.
\newblock Tuning-free stochastic optimization.
\newblock In \emph{Proceedings of the 41st International Conference on
Machine Learning}, pages 23622--23661, 2024.
\newblock \url{https://proceedings.mlr.press/v235/khaled24a.html}.

\bibitem[Kingma and Ba(2015)]{kingma2015}
Diederik P. Kingma and Jimmy Ba.
\newblock Adam: A method for stochastic optimization.
\newblock In \emph{International Conference on Learning Representations},
2015.
\newblock \url{https://arxiv.org/abs/1412.6980}.

\bibitem[Li et al.(2025)]{li2025adamw}
Huan Li, Yiming Dong, and Zhouchen Lin.
\newblock On the $O(\sqrt{d}/K^{1/4})$ convergence rate of AdamW measured by
$\ell_1$ norm.
\newblock 2025.
\newblock \url{https://arxiv.org/abs/2505.11840}.

\bibitem[Lessard et al.(2016)]{lessard2016}
Laurent Lessard, Benjamin Recht, and Andrew Packard.
\newblock Analysis and design of optimization algorithms via integral
quadratic constraints.
\newblock \emph{SIAM Journal on Optimization}, 26(1):57--95, 2016.
\newblock \url{https://doi.org/10.1137/15M1009597}.

\bibitem[Loshchilov and Hutter(2017)]{loshchilov2017sgdr}
Ilya Loshchilov and Frank Hutter.
\newblock SGDR: Stochastic gradient descent with warm restarts.
\newblock In \emph{International Conference on Learning Representations},
2017.
\newblock \url{https://arxiv.org/abs/1608.03983}.

\bibitem[Loshchilov and Hutter(2019)]{loshchilov2019}
Ilya Loshchilov and Frank Hutter.
\newblock Decoupled weight decay regularization.
\newblock In \emph{International Conference on Learning Representations},
2019.
\newblock \url{https://openreview.net/forum?id=Bkg6RiCqY7}.

\bibitem[Nesterov(1983)]{nesterov1983}
Yurii E. Nesterov.
\newblock A method for solving the convex programming problem with convergence
rate $O(1/k^2)$.
\newblock \emph{Doklady Akademii Nauk SSSR}, 269:543--547, 1983.

\bibitem[Nesterov(2004)]{nesterov2004}
Yurii Nesterov.
\newblock \emph{Introductory Lectures on Convex Optimization: A Basic Course}.
\newblock Kluwer Academic Publishers, 2004.
\newblock \url{https://doi.org/10.1007/978-1-4419-8853-9}.

\bibitem[Nesterov(2015)]{nesterov2015}
Yurii Nesterov.
\newblock Universal gradient methods for convex optimization problems.
\newblock \emph{Mathematical Programming}, 152:381--404, 2015.
\newblock \url{https://doi.org/10.1007/s10107-014-0790-0}.

\bibitem[O'Donoghue and Cand\`es(2015)]{odonoghue2015}
Brendan O'Donoghue and Emmanuel J. Cand\`es.
\newblock Adaptive restart for accelerated gradient schemes.
\newblock \emph{Foundations of Computational Mathematics}, 15(3):715--732,
2015.
\newblock \url{https://doi.org/10.1007/s10208-013-9150-3}.

\bibitem[Polyak(1964)]{polyak1964}
Boris T. Polyak.
\newblock Some methods of speeding up the convergence of iteration methods.
\newblock \emph{USSR Computational Mathematics and Mathematical Physics},
4(5):1--17, 1964.
\newblock \url{https://doi.org/10.1016/0041-5553(64)90137-5}.

\bibitem[Reddi et al.(2018)]{reddi2018}
Sashank J. Reddi, Satyen Kale, and Sanjiv Kumar.
\newblock On the convergence of Adam and beyond.
\newblock In \emph{International Conference on Learning Representations},
2018.
\newblock \url{https://openreview.net/forum?id=ryQu7f-RZ}.

\bibitem[Roulet and d'Aspremont(2020)]{roulet2020}
Vincent Roulet and Alexandre d'Aspremont.
\newblock Sharpness, restart and acceleration.
\newblock \emph{SIAM Journal on Optimization}, 30(1):262--289, 2020.
\newblock \url{https://arxiv.org/abs/1702.03828}.

\bibitem[Soudry et al.(2018)]{soudry2018}
Daniel Soudry, Elad Hoffer, Mor Shpigel Nacson, Suriya Gunasekar, and Nathan
Srebro.
\newblock The implicit bias of gradient descent on separable data.
\newblock \emph{Journal of Machine Learning Research}, 19(70):1--57, 2018.
\newblock \url{https://www.jmlr.org/papers/v19/18-188.html}.

\bibitem[Sutskever et al.(2013)]{sutskever2013}
Ilya Sutskever, James Martens, George Dahl, and Geoffrey Hinton.
\newblock On the importance of initialization and momentum in deep learning.
\newblock In \emph{Proceedings of the 30th International Conference on
Machine Learning}, pages 1139--1147, 2013.
\newblock \url{https://proceedings.mlr.press/v28/sutskever13.html}.

\bibitem[Touvron et al.(2023)]{touvron2023llama2}
Hugo Touvron et al.
\newblock Llama 2: Open foundation and fine-tuned chat models.
\newblock 2023.
\newblock \url{https://arxiv.org/abs/2307.09288}.

\bibitem[Vaswani et al.(2026)]{vaswani2026}
Sharan Vaswani, Yifan Sun, and Reza Babanezhad.
\newblock Convergence of steepest descent and Adam under non-uniform
smoothness.
\newblock 2026.
\newblock \url{https://arxiv.org/abs/2605.30648}.

\bibitem[Wang et al.(2024)]{wang2024}
Bohan Wang, Yushun Zhang, Huishuai Zhang, Qi Meng, Ruoyu Sun, Zhi-Ming Ma,
Tie-Yan Liu, Zhi-Quan Luo, and Wei Chen.
\newblock Provable adaptivity of Adam under non-uniform smoothness.
\newblock In \emph{Proceedings of the 30th ACM SIGKDD Conference on Knowledge
Discovery and Data Mining}, 2024.
\newblock \url{https://arxiv.org/abs/2208.09900}.

\bibitem[Xie et al.(2025)]{xie2025}
Shuo Xie, Mohamad Amin Mohamadi, and Zhiyuan Li.
\newblock Adam exploits $\ell_\infty$ geometry of loss landscape via
coordinate-wise adaptivity.
\newblock In \emph{International Conference on Learning Representations},
2025.
\newblock \url{https://proceedings.iclr.cc/paper_files/paper/2025/file/1700ad4e6252e8f2955909f96367b34d-Paper-Conference.pdf}.

\bibitem[Yue et al.(2023)]{yue2023}
Pengyun Yue, Cong Fang, and Zhouchen Lin.
\newblock On the lower bound of minimizing Polyak--\L{}ojasiewicz functions.
\newblock 2023.
\newblock \url{https://arxiv.org/abs/2212.13551}.

\bibitem[Zhang et al.(2019)]{zhang2019ngd}
Guodong Zhang, James Martens, and Roger B. Grosse.
\newblock Fast convergence of natural gradient descent for over-parameterized
neural networks.
\newblock In \emph{Advances in Neural Information Processing Systems 32},
2019.
\newblock \url{https://proceedings.neurips.cc/paper_files/paper/2019/hash/1da546f25222c1ee710cf7e2f7a3ff0c-Abstract.html}.

\bibitem[Zhang et al.(2024a)]{zhang2024adambias}
Chenyang Zhang, Difan Zou, and Yuan Cao.
\newblock The implicit bias of Adam on separable data.
\newblock In \emph{Advances in Neural Information Processing Systems 37},
2024.
\newblock \url{https://arxiv.org/abs/2406.10650}.

\bibitem[Zhang et al.(2024b)]{zhang2024transformers}
Yushun Zhang, Congliang Chen, Tian Ding, Ziniu Li, Ruoyu Sun, and Zhi-Quan Luo.
\newblock Why transformers need Adam: A Hessian perspective.
\newblock In \emph{Advances in Neural Information Processing Systems 37},
2024.
\newblock \url{https://proceedings.neurips.cc/paper_files/paper/2024/hash/ee0e45ff4de76cbfdf07015a7839f339-Abstract-Conference.html}.

\bibitem[Zhuang et al.(2022)]{zhuang2022}
Zhenxun Zhuang, Mingrui Liu, Ashok Cutkosky, and Francesco Orabona.
\newblock Understanding AdamW through proximal methods and scale-freeness.
\newblock \emph{Transactions on Machine Learning Research}, 2022.
\newblock \url{https://arxiv.org/abs/2202.00089}.

\bibitem[Zou et al.(2019)]{zou2019}
Fangyu Zou, Li Shen, Zequn Jie, Weizhong Zhang, and Wei Liu.
\newblock A sufficient condition for convergences of Adam and RMSProp.
\newblock In \emph{Proceedings of the IEEE/CVF Conference on Computer Vision
and Pattern Recognition}, pages 11127--11135, 2019.
\newblock \url{https://arxiv.org/abs/1811.09358}.

\end{thebibliography}

\end{document}
